% SOURCE docs/release/2026-10-29/source/docs/writing-repair/2026-10-29/src/f09708095a17fce9.tex CHAR 0
\chapter{Effect estimation, heterogeneous actions and threshold comparisons}
\label{exec17:s4}

A servicing team may want to know the value of offering assistance to every
eligible account, the value among accounts already receiving assistance, or
the value in the population where receipt remains uncertain. Those questions
need not have the same answer. Nor does learning a large conditional benefit
identify the best use of a scarce specialist team. A threshold experiment can
support a local comparison while leaving the portfolio decision unresolved.
The proposal therefore needs a vocabulary for comparing estimators through
their population, uncertainty and decision consequences.

This chapter develops three connections: assignment-dependent populations and
orthogonal estimation; joint heterogeneous comparisons and resource allocation;
and threshold identification versus finite approximation. The numerical laws
are stipulated, reproducible examples in USD per eligible account over a fixed
thirty-day horizon. They establish algebra and decision consequences under
those laws. They are neither bank observations nor new executions of the
external software discussed below.~\citep{exec17_s4_math}

The later comparisons follow these distinctions into individual software
paths: the population actually averaged, the provenance of nuisance predictions,
and the construction of the final uncertainty. Repeated examples use the same
stipulated customers where that makes a difference between implementations
visible. The alternatives remain separate when their priors, sampling units,
missingness rules or numerical conventions answer different questions.

\section{The population can depend on the assignment mechanism}
\label{exec17:s4-target}

Let $O=(X,A,Y)$ denote one independent account observation, where $X$ is
information available before a binary action $A$, and $Y$ is the subsequent
monetary outcome. Assume consistency, no interference, conditional
exchangeability of the two potential outcomes given $X$, and positive
probabilities for both actions on the population being studied. These are
scientific restrictions; fitting a flexible regression does not verify them.
Define
\[
 e(x)=P(A=1\mid X=x),\qquad m_a(x)=E(Y\mid A=a,X=x),\qquad
 \tau(x)=m_1(x)-m_0(x).
\]
Throughout the following derivations, the moments exist and the denominators
are nonzero. Uniform overlap and suitable higher moments will be imposed when
an asymptotic variance or remainder bound needs them.

A useful family of targets is
\begin{equation}
 \theta_h(P)=\frac{E\{h(e(X))\tau(X)\}}{H},\qquad
 H=E\{h(e(X))\}>0.
 \label{exec17:s4-weight-target}
\end{equation}
The cases $h=1$, $h(e)=e$, $h(e)=1-e$ and $h(e)=e(1-e)$ give the population,
treated, untreated and overlap-weighted effects. The treated case follows
because $E\{e(X)f(X)\}=P(A=1)E\{f(X)\mid A=1\}$. Overlap weighting emphasizes
accounts with uncertain receipt; it does not estimate the value of treating
every account. The bank must decide whether that emphasis describes an
actionable population or only a useful scientific comparison.

There is a further distinction within the word ``weighted.'' A weight
$w(x)$ supplied independently and held fixed defines
$E\{w(X)\tau(X)\}/Ew(X)$. A weight computed from the unknown assignment law
defines a functional that changes when that law changes. Treating an estimated
overlap weight as fixed in a variance calculation can discard a component of
uncertainty. A randomized experiment with a known, fixed assignment design is
yet another model: its allowed perturbations do not include arbitrary changes
to that assignment law. The relevant derivative must match the model being
claimed.~\citep{exec17_s4_math,dml2018}

To see the missing component, consider a smooth parametric path through the
observable law, with mean-zero score $s(O)$. A dot denotes a path derivative.
Conditional differentiation gives
\begin{align*}
 \dot e(x)&=E\{(A-e(x))s(O)\mid X=x\},\\
 \dot m_a(x)&=E\{(Y-m_a(x))s(O)\mid A=a,X=x\}.
\end{align*}
The marginal distribution of $X$ also moves. Applying the quotient rule to
Equation~\eqref{exec17:s4-weight-target} separates that movement from changes
in $m_a$ and $e$. For differentiable $h$, a mean-zero influence function in
the unrestricted observational model is
\begin{align}
 \phi_h(O)={}&\frac{h(e)}{H}
 \left\{\tau-\theta_h+\frac{A}{e}(Y-m_1)
                 -\frac{1-A}{1-e}(Y-m_0)\right\}\notag\\
 &+\frac{h'(e)}{H}(A-e)(\tau-\theta_h).
 \label{exec17:s4-weight-if}
\end{align}
Functions of $X$ in this expression are evaluated at the observed $X$.
Taking its covariance with $s$ recovers the three derivative contributions.
The last line is the assignment-law contribution. It disappears for fixed
external weights, but generally survives for treated and overlap targets.
This is a project derivation under the stated model, rather than an inference
that any package implements every member of this family.~\citep{exec17_s4_math}

The extra term is orthogonal to the first line. Its conditional mean given
$X$ is zero; its covariance with each outcome residual is zero because
$E(Y-m_a\mid A=a,X)=0$. Consequently its addition increases the variance by
\begin{equation}
 \frac{E\{h'(e)^2e(1-e)(\tau-\theta_h)^2\}}{H^2}.
 \label{exec17:s4-extra-variance}
\end{equation}
This comparison concerns the two displayed random variables in this
observational model. It does not claim that every estimator using a different
model must have that variance, or that weights learned on an independent
sample have the same inferential target.

For $h(e)=e$, write $p=P(A=1)$. Combining terms in
Equation~\eqref{exec17:s4-weight-if} yields
\begin{equation}
 \phi_{\rm ATT}=
 \frac{A}{p}(Y-m_0-\theta_{\rm ATT})
 -\frac{(1-A)e}{p(1-e)}(Y-m_0).
 \label{exec17:s4-att-if}
\end{equation}
The treated-arm outcome regression cancels. This agrees with the distinct
ATE and ATTE scores developed by \citet{dml2018}; it explains why an adapter
must preserve the target-specific score and nuisance requirements. Simply
relabeling an ATE result as ATT, or retaining ATE nuisance diagnostics as if
they were necessary and sufficient for ATT, changes the mathematical claim.

Hard trimming requires additional care. With
$h(e)=1\{\ell\leq e\leq u\}$, the boundary can move under a path, and the
ordinary smooth derivative above is not automatically available. A declared
fixed subpopulation, a smooth weight, and a population selected by estimated
propensity thresholds are three different specifications. Mass near the
thresholds, whether selection uses an independent sample, and whether the
estimand conditions on the selected set matter. Removing unstable accounts
from a calculation does not by itself justify an interval for the original
population.

\section{An overlap score and its exact remainder}
\label{exec17:s4-remainder}

Orthogonality is a local derivative statement. It is useful precisely because
it can suppress first-order nuisance error, but that statement must not be
silently strengthened into consistency under arbitrary misspecification.
Consider the overlap case $h(e)=e(1-e)$. Let $q(x)$ be a fitted propensity and
$\widehat m_a(x)$ fitted arm means, evaluated on observations excluded from
their training folds. Define
\begin{align}
 D(O;q)&=A(1-q)^2+(1-A)q^2,\notag\\
 N(O;q,\widehat m)&=D(O;q)(\widehat m_1-\widehat m_0)
      +A(1-q)(Y-\widehat m_1)
      -(1-A)q(Y-\widehat m_0),\notag\\
 \widehat\theta&=\frac{\mathbb P_nN}{\mathbb P_nD}.
 \label{exec17:s4-overlap-ratio}
\end{align}
At the truth, $E D=H$. Also
$D=e(1-e)+(1-2e)(A-e)$ for binary $A$, so the centered ratio score
$(N-\theta_hD)/H$ equals Equation~\eqref{exec17:s4-weight-if}.
The assignment derivative has been incorporated into both numerator and
denominator; it is not an optional standard-error correction applied after
estimating a different target.

The bias can be computed without an asymptotic approximation. Put
$\delta=q-e$ and $\delta_a=\widehat m_a-m_a$. Conditional expectation gives
\begin{align}
 E(D\mid X)&=e(1-e)+\delta^2,\notag\\
 E(N-\theta_hD)&=
 E\left[\delta^2(\tau-\theta_h)
       +\delta\{q\delta_1+(1-q)\delta_0\}\right].
 \label{exec17:s4-overlap-remainder}
\end{align}
For example, the coefficient on $\delta_1$ before simplification is
$E(D\mid X)-e(1-q)=q(q-e)$; the coefficient on $\delta_0$ is
$-E(D\mid X)+(1-e)q=(1-q)(q-e)$. These identities explain both the
second-order propensity term and the product term. The population fitted
ratio differs from $\theta_h$ by the second line divided by
$H+E\delta^2$.~\citep{exec17_s4_math}

If $q=e$, the remainder vanishes even with inaccurate arm means. If the arm
means are correct but $q$ is arbitrary, the first term generally remains.
Thus this corrected overlap construction is locally orthogonal while lacking
global double robustness in that second direction. For a fixed external
weight and an ATE-style augmented score, the usual product remainder is a
different calculation. The word ``augmented'' cannot resolve the distinction.

Under bounded $\tau-\theta_h$ and fitted probabilities in $[0,1]$, a bound
on the absolute numerator error is
\[
 \|\tau-\theta_h\|_\infty\|q-e\|_2^2+
 \|q-e\|_2\{\|\widehat m_1-m_1\|_2+
                          \|\widehat m_0-m_0\|_2\}.
\]
Sufficient conditions for this bound to be $o_p(n^{-1/2})$ include a
squared-propensity-error condition and a propensity/outcome product-rate
condition, with a denominator bounded away from zero. An asymptotic linear
expansion also needs control of the fitted score in $L_2$, moments, and the
fold evaluation terms. Cross-fitting addresses training/evaluation dependence;
it neither proves those rates nor supplies exchangeability. The conditions in
the DML paper similarly accompany its own specified scores and statistical
model. They are not a theorem for an arbitrary callable score.~\citep{dml2018,doubleml2022}

In a bank proposal these conditions create a useful design decision. An
overlap target can avoid assigning large inverse weights to nearly deterministic
receipt, but the target itself uses the receipt mechanism. An independently
specified segment can have a simpler fixed-weight target while retaining
difficult support. Neither option makes the scientific population choice
unnecessary. The appropriate comparison holds the target fixed first and
then studies estimation error, rather than treating every numerically stable
answer as an improvement.

\section{Clipping, population change and apparent precision}
\label{exec17:s4-clipping}

Consider three pre-action strata with the following exact law:
\[
\begin{array}{c|rrrr}
 X& P(X)&e(X)&m_0(X)\ ({\rm USD})&\tau(X)\ ({\rm USD})\\ \hline
 1&1/4&1/10&30&-4\\
 2&1/2&1/2&50&4\\
 3&1/4&9/10&70&20
\end{array}
\]
Within each stratum the action is generated independently of potential
outcomes. The ATE is $6$, ATT is $54/5$, and ATC is $6/5$ USD. Those basic
target differences motivate the example; the new issue is what fitted
weight changes do to its score and uncertainty. The true overlap mass is
$H=17/100$, giving $\theta_h=86/17$ USD. True-propensity trimming to
$[1/5,4/5]$ retains only the middle stratum and gives $4$ USD.
\citep{exec17_s4_math}

Now clip the two extreme propensities to obtain
$q=(1/5,1/2,4/5)$. The unnormalized plain inverse-weighted contrast
\[
 E\left\{\frac{AY}{q}-\frac{(1-A)Y}{1-q}\right\}
\]
equals $107/8$ USD, despite an ATE of $6$. Its coefficients on baseline
outcome levels no longer cancel in each stratum. This is a population
calculation, so more observations do not remove the discrepancy. An ATE
augmented score with correct arm means still gives $6$; its residual terms
have conditional mean zero. That result does not apply to every weighted
target simply because the same fitted $q$ is used.

If one instead freezes the weights $q(1-q)$ and averages the true effects,
the target is $228/41$ USD. The corrected overlap ratio in
Equation~\eqref{exec17:s4-overlap-ratio}, still using correct arm means,
gives $36/7$ USD. Its bias relative to the true overlap effect is $10/119$.
All four quantities---ATE, true overlap, frozen fitted-weight average and
corrected fitted ratio---can be estimated precisely while answering different
questions. With the stipulated arm-mean errors
$\delta_1=(1,-2,3)$ and $\delta_0=(-1,1,2)$ USD, the corrected ratio's
unnormalized error is $-239/3400$ USD and its denominator is $7/40$;
the normalized error is $-239/595$ USD. The exact decomposition therefore
checks the sign and size of a genuine product remainder, rather than merely
showing that a program terminates.~\citep{exec17_s4_math}

The familiar effective-sample-size diagnostic can also look better after
clipping. For the treated weight $W=A/q(X)$, its population analogue is
\[
 \frac{\operatorname{ESS}}{n}\longrightarrow
       \frac{\{E(W)\}^2}{E(W^2)}.
\]
Under the true propensity it is $9/34$; under the clipped propensity it is
$841/2024$. The larger second fraction measures concentration of those
weights. It neither proves target preservation nor yields an outcome-specific
standard error: outcome noise, effect heterogeneity and fitted-score error
still matter. A self-normalized estimator would be a further algorithm with
its own population ratio, not a justification for the unnormalized contrast
used above.

The assignment derivative can be checked directly in this same law. Perturb
the three propensities by $t(1/20,-1/10,1/20)$, holding stratum masses and
outcome laws fixed. The ATT derivative is $27/50$ USD per unit $t$ and the
overlap-effect derivative is $-24/17$. The covariance of the first line of
Equation~\eqref{exec17:s4-weight-if} alone with the assignment score is zero;
the complete influence function gives exactly these derivatives. Mean zero
of a candidate influence function is therefore insufficient. The derivative
test detects what a simple centering test misses.~\citep{exec17_s4_math}

\section{Selected software paths embody different estimators}
\label{exec17:s4-paths}

The inspected DoubleML paper describes linear score architecture and
model-specific nuisance estimation, with extension routes for callable scores
and sample splitting. The inherited local execution used the selected IRM
ATE path in version 0.11.4, released 10 August 2026, with externally supplied
nuisances and folds, no IPW normalization and truncation at 0.01. That is
evidence for its recorded operation only. The inspected source adjusts fitted
propensities before score evaluation. Changing truncation, target or callable
score would require a new comparison, including the remainder appropriate to
that functional.~\citep{doubleml2022,runtime_doubleml_scope,exec17_s4_paths}

The pinned R \texttt{grf} implementation takes a different route for its
binary-treatment overlap target. Its average-effect source fits a regression
of outcome residuals on treatment residuals, with an intercept, and uses a
heteroskedastic or cluster sandwich calculation. It does not simply run the
all-population AIPW average with another averaging weight. At true nuisance
functions, write $\mu(X)=E(Y\mid X)$ and
$Y-\mu=(A-e)\tau+\varepsilon$, with $E(\varepsilon\mid A,X)=0$.
The residual-regression slope is
\[
 \frac{E\{(A-e)(Y-\mu)\}}{E\{(A-e)^2\}}
       =\frac{E\{e(1-e)\tau\}}{H}.
\]
Because both residuals have mean zero at truth, an intercept does not change
this population slope. Its first-order residual-regression influence function
at truth is
\[
 \frac{(A-e)^2(\tau-\theta_h)+(A-e)\varepsilon}{H},
\]
which equals the full overlap influence function in
Equation~\eqref{exec17:s4-weight-if} for binary $A$.
This algebra connects the selected implementation to the assignment-dependent
target.~\citep{exec13_grf_rpath,exec17_s4_math}

Equality at truth does not make fitted algorithms identical. In the clipped
three-stratum example, retain the true marginal outcome mean $\mu$ but use
$q$ for treatment residuals. Here $E(A-q)=0$,
$E(q-e)^2=1/200$, and the residual-regression slope is
\[
 \frac{E\{e(1-e)\tau\}}{H+E(q-e)^2}=\frac{172}{35}.
\]
It differs both from the true overlap effect $86/17$ and from the corrected
ratio $36/7$ using correct arm means. These nuisance specifications are
deliberately different: true $\mu$ is not the same fitted object as the two
true arm means combined with an incorrect $q$. The example shows why a
common target and common limiting influence function do not establish
finite fitted equivalence.~\citep{exec17_s4_math}

At revision \texttt{5bee99b5}, the selected \texttt{grf} source's overlap branch
does not use the method choice after entering that branch. Earlier guards
still reject particular TMLE calls with clustering or sample weights. Thus
``method ignored'' must be qualified by control flow. Other inspected forest
types also have restrictions on target and method combinations. An adapter
should validate the exact combination, preserve package diagnostics and
record the actual route used. The generalized-forest paper's pointwise
normality result is conditional on its assumptions and honest, regular,
subsampled forest construction; it does not establish inference for every
selected policy or every fitted forest.~\citep{exec13_grf_rpath,exec13_grf_native,grf2018}

These comparisons favor wrapping selected public interfaces before porting
their internals. Reproducing an R or Python score in a common C++ core could
eventually reduce interprocess overhead, but would create an additional
statistical implementation to maintain. First establish the mathematical
target, nuisance convention, supported uncertainty unit and failure behavior
on matched examples. Reject unsupported target/method combinations at that
boundary. No current source inspection establishes complete dependency rights
or the performance of a new wrapper.~\citep{exec17_s4_design}

\section{Shared controls make heterogeneous comparisons joint}
\label{exec17:s4-joint}

Suppose the bank can choose no contact, a message or specialist assistance,
indexed by $A=0,1,2$. Let $e_j(X)$ and $m_j(X)$ denote action probabilities
and arm means. Under the corresponding multiarm consistency, exchangeability
and support assumptions, define every contrast on the same eligible-account
law:
\[
 \theta_j=E\{m_j(X)-m_0(X)\},\qquad j=1,2.
\]
Its population influence function is
\[
 \phi_j=\tau_j(X)-\theta_j+
       \frac{1\{A=j\}}{e_j}(Y-m_j)
       -\frac{1\{A=0\}}{e_0}(Y-m_0).
\]
For distinct noncontrol actions $j,k$, mutually exclusive arm indicators
eliminate their outcome-residual cross-products, but the shared control
residual appears twice. Therefore
\begin{equation}
 \operatorname{Cov}(\phi_j,\phi_k)=
 \operatorname{Cov}\{\tau_j(X),\tau_k(X)\}
       +E\{\sigma_0^2(X)/e_0(X)\},
 \label{exec17:s4-shared-control}
\end{equation}
where $\sigma_0^2(X)=\operatorname{Var}(Y\mid A=0,X)$.
Conditional mean-zero residuals remove the remaining cross terms. The first
term can have either sign; a positive shared-control term does not imply
positive total covariance in every heterogeneous population.

Take constant effects, assignment probabilities $(1/2,1/4,1/4)$, and
conditional outcome-noise variance $4$ USD$^2$ in all arms. The influence
covariance matrix and variance of the difference are
\[
 \Sigma=\begin{pmatrix}24&8\\8&24\end{pmatrix},\qquad
 \operatorname{Var}(\phi_2-\phi_1)=32.
\]
Adding marginal variances as though the contrasts were independent would give
$48$. In this example the shared control cancels in the difference. For an
iid mean, the corresponding asymptotic covariance is $\Sigma/n$; these are
not empirical standard errors and are not formulas for clustered bank
histories. Joint confidence regions or multiple-comparison decisions require
an explicitly chosen simultaneous procedure. Two marginal intervals alone
do not specify one.~\citep{exec17_s4_math}

This issue is concrete in the pinned CausalML T-learner path. Its point mean
uses the full prediction covariate matrix, while the inspected analytic
variance uses the observations in the current treatment/control pair. Under
unequal assignment or covariate composition, those sets can represent
different populations. The finding is about the selected revision and path,
not a claim that every CausalML estimator behaves that way. A proposed
multiarm wrapper should standardize every action prediction to one explicit
evaluation population, then construct the joint score and covariance there.
It should keep package-specific uncertainty as a separately named result
until equivalence is established.~\citep{exec14_causalml_primary,exec17_s4_paths}

The same source offers a row bootstrap. Rows from repeated accounts or a
shared calendar shock are not independent sampling units simply because they
have distinct row indices. Whether the comparator uses independent accounts,
households, branches or time blocks belongs in the statistical specification.
A bootstrap setting cannot repair an undeclared unit or create joint support
for an action rarely observed in a subgroup.

\section{A resource constraint changes the optimal action}
\label{exec17:s4-allocation}

Let $g(x)=\tau(x)-c(x)$ be the incremental net value of one action, where
$c(x)$ includes the action's incremental cost in the same USD units and
horizon. Let $r(x)>0$ be specialist resource units and let $d(x)\in[0,1]$
be the fraction of that stratum assigned to the action. With expected
resource budget $B$ per eligible account, consider
\begin{equation}
 \max_d E\{g(X)d(X)\}\quad\text{subject to}\quad
 E\{r(X)d(X)\}\leq B,\quad 0\leq d\leq1.
 \label{exec17:s4-resource-problem}
\end{equation}
Fractional decisions describe randomization or divisible allocation in this
population problem. They are not a guarantee that a finite day's integer
roster fits a hard capacity constraint.

For a multiplier $\lambda\geq0$, every feasible policy obeys the dual bound
\begin{equation}
 E(gd)\leq \lambda B+E\{(g-\lambda r)_+\}.
 \label{exec17:s4-resource-dual}
\end{equation}
Indeed, $gd=(g-\lambda r)d+\lambda rd$ and
$(g-\lambda r)d\leq(g-\lambda r)_+$. Equality is attained if the policy
uses the action wherever $g>\lambda r$, avoids it wherever $g<\lambda r$,
and allocates any equality stratum to meet the budget when the multiplier is
positive. This is a checkable optimality certificate. It explains why ranking
by effect, or even by net gain without resource units, can be wrong.

Consider four strata:
\[
\begin{array}{c|rrrr}
 \text{stratum}&1&2&3&4\\ \hline
 P(X)&2/5&3/10&1/5&1/10\\
 \text{gross effect, USD}&4&9&15&25\\
 \text{incremental cost, USD}&1&3&7&14\\
 g(X),\ \text{USD}&3&6&8&11\\
 r(X),\ \text{units}&1&1&2&4
\end{array}
\]
Set $B=1/2$ unit per eligible account. The policy
$d=(0,1,1/2,0)$ uses exactly $1/2$ unit and produces $13/5$ USD of net value.
At $\lambda=4$ USD per resource unit, the adjusted gains
$g-\lambda r$ are $(-1,2,0,-5)$; the dual bound is
$4(1/2)+(3/10)2=13/5$. The feasible allocation therefore attains the global
bound. Sorting by gross effect or by unadjusted net gain fills stratum 4
first and then a quarter of stratum 3, giving
$d=(0,0,1/4,1)$ and $3/2$ USD. The loss is $11/10$ USD per eligible
account, even though the ranked effects are known exactly.~\citep{exec17_s4_math}

This example deliberately separates scientific estimation from operational
optimization. The cost schedule is stipulated, specialist time is a separate
constraint, and the value is incremental relative to no action. If monetary
cost already prices the same scarce time, the sponsor must explain whether
the hard constraint remains necessary and avoid counting the same resource
twice. Risk limits, fairness conditions and action eligibility can further
restrict the feasible set. For indivisible actions on a finite roster, the
appropriate optimization can become a knapsack or assignment problem; the
fractional certificate does not assert that a greedy integer solution is
optimal.

For multiple actions, replace $d$ by $d_j(x)\geq0$ with
$\sum_jd_j(x)\leq1$, action-specific gains and resources. At a fixed
multiplier, each account receives the action with the largest positive
$g_j-\lambda r_j$, subject to any additional declared constraints. Shared
controls and common target populations remain essential because the optimizer
compares these action values directly. Exporting unrelated pairwise targets
to an optimizer can produce a feasible but scientifically meaningless plan.

\section{From effect error to decision loss}
\label{exec17:s4-regret}

Prediction error and policy error need not move together. For the same fixed
feasible policy class $\mathcal D$, let $d^*$ maximize $E(gd)$ and
$\widehat d$ maximize $E(\widehat g d)$. Suppose the pointwise bound
$|\widehat g-g|\leq\epsilon(X)$ is valid. Adding and subtracting the fitted
objective gives
\begin{align}
 E\{g(d^*-\widehat d)\}
 &=E\{(g-\widehat g)(d^*-\widehat d)\}
       +E\{\widehat g(d^*-\widehat d)\}\notag\\
 &\leq E\{\epsilon(X)|d^*-\widehat d|\}.
 \label{exec17:s4-regret-bound}
\end{align}
The last fitted-objective difference is nonpositive by optimality. If the
optimizer is within $\eta$ USD of its optimum, add $\eta$ to the bound.
This is a deterministic comparison conditional on the two gain functions,
the common population objective and the same feasible class. It is not a
confidence bound supplied by an arbitrary learner. If the fitted optimizer
uses an empirical objective, a population-to-empirical discrepancy must also
be controlled. If costs or capacity are estimated, their errors can change
the feasible set, so this derivation no longer covers them by itself.

With $0\leq d\leq1$, a uniform gain-error bound $\epsilon_0$ gives the coarse
bound $\epsilon_0$ per account. The sharper expression charges error only
where the two policies disagree. This explains why errors near a resource
threshold can be more consequential than equally large errors in strata
where the same action remains optimal. A low global squared prediction loss
does not alone demonstrate low constrained decision loss. The finite
reference enumerates 81 signed-error configurations over the binary feasible
policies and confirms the inequality, while the algebra establishes its stated
general domain.~\citep{exec17_s4_math}

The observational policy-learning analysis of \citet{policy2021} supplies a
more substantive statistical route: construct doubly robust scores and learn
within a controlled-complexity policy class under overlap, noise and nuisance
rate conditions. Its regret theorem is tied to the admitted binary-decision
problem and assumptions. It does not automatically cover every resource
constraint, multiarm action, dependent history or optimizer introduced here.
Likewise, a pointwise conditional-effect interval from a forest is not an
interval for the value of a policy selected using the same sample.

A proposal comparison should therefore split learning and evaluation at the
appropriate independent unit, declare the reference population and feasible
class, and evaluate the frozen selected policy with a justified value score.
It should report infeasibility and inability to estimate value as outcomes.
An exact reference can verify the optimizer; a prospective generated-outcome
study can investigate selection and finite-sample inference. Neither supplies
the true net effect of an untested bank action.

\section{Threshold estimates include an approximation problem}
\label{exec17:s4-curvature}

Now let $S$ be a pre-action eligibility score, centered at a cutoff of zero.
A sharp rule assigns the action when $S\geq0$. Continuity of the relevant
potential-outcome conditional means at zero, suitable support approaching
the cutoff and a well-defined assignment rule can identify a local cutoff
effect. Sorting, manipulation, concurrent threshold rules and delivery
versions can invalidate that interpretation. The following calculation
assumes the scientific cutoff question has already been justified and asks
what a finite-bandwidth approximation estimates.

On the right side write $S=hu$, with $u\in[0,1]$, and take uniform local
design density and uniform weights for an exact population local-linear
projection. Suppose the side's conditional mean is
$b_0+b_1S+\kappa S^2$. The normal equations are
\[
 G\begin{pmatrix}\alpha\\\beta\end{pmatrix}
 =\int_0^1\begin{pmatrix}1\\u\end{pmatrix}
       (b_0+b_1hu+\kappa h^2u^2)\,du,
 \qquad
 G=\begin{pmatrix}1&1/2\\1/2&1/3\end{pmatrix}.
\]
Since $G^{-1}=\left(\begin{smallmatrix}4&-6\\-6&12\end{smallmatrix}\right)$,
the intercept equals
\begin{equation}
 \alpha=\int_0^1(4-6u)m(hu)\,du
        =b_0-\frac{\kappa h^2}{6}.
 \label{exec17:s4-local-bias}
\end{equation}
The equivalent kernel integrates to one and has zero first moment, but its
second moment is $-1/6$. Its negative weights beyond $u=2/3$ are part of
boundary trend correction. They are not positive population weights.
Here $\kappa$ is the quadratic coefficient; the second derivative is
$2\kappa$. Confusing the two would double the bias correction.
\citep{exec17_s4_math}

Using $S=-hu$ on the left changes the linear term's sign and leaves the
quadratic intercept bias unchanged. With equal bandwidths, the difference
between projected intercepts has bias
$-(\kappa_+-\kappa_-)h^2/6$. Unequal bandwidths replace this by
$-(\kappa_+h_+^2-\kappa_-h_-^2)/6$. These are exact projection identities
for the stipulated quadratic and uniform-density law. A general density,
another kernel or a nonquadratic regression requires its own expansion and
remainder. A finite random-sample local regression is not identical to the
population projection; its expected ratio of random matrices is another
object.

Reducing bandwidth can reduce approximation bias while removing nearby
observations and destabilizing the local design. Estimated curvature,
bandwidth selection and robust bias correction add their own uncertainty.
An interval using a variance formula for an uncorrected intercept cannot be
declared valid merely because a point bias adjustment was subtracted. The
inspected \texttt{rdrobust} help documents local-polynomial, bias-corrected
and variance choices; the formula above is a transparent reference for a
particular population calculation, not a substitute for its full inference
theory.~\citep{exec13_ivrd_rdrobust}

\section{A fuzzy ratio can amplify approximation error}
\label{exec17:s4-fuzzy}

With incomplete receipt, let $\Delta_Y$ and $\Delta_A$ denote the true
outcome and receipt discontinuities. Under the separately required exclusion,
monotonicity and local continuity restrictions, their ratio
$\theta=\Delta_Y/\Delta_A$ has a local complier interpretation when the
receipt jump is nonzero. Let $b_Y,b_A$ be the biases in the population
projected discontinuities. Exact algebra gives
\begin{equation}
 \frac{\Delta_Y+b_Y}{\Delta_A+b_A}-\theta
       =\frac{b_Y-\theta b_A}{\Delta_A+b_A}.
 \label{exec17:s4-fuzzy-bias}
\end{equation}
Even before sampling uncertainty, an apparently small bias in the receipt
jump can change the denominator's magnitude or sign. Correcting only the
outcome projection does not generally correct the ratio.

Set the true local effect to $8$ USD, the outcome curvature-coefficient
difference to $12$ USD per score-unit squared and the receipt difference to
$1/2$ per score-unit squared. Then $b_Y=-2h^2$ USD and $b_A=-h^2/12$.
The exact projected ratios are
\[
\begin{array}{c|cc|c}
 \Delta_A&h&\Delta_A+b_A&
    (\Delta_Y+b_Y)/(\Delta_A+b_A)\ ({\rm USD})\\ \hline
 1/2&1/2&23/48&168/23\\
 1/2&1/4&95/192&744/95\\
 1/50&1/2&-1/1200&408\\
 1/50&1/4&71/4800&168/71
\end{array}
\]
All rows have the same true ratio of $8$. At the weak receipt jump, the
first projected ratio is $408$ USD, although no sampling noise has been
introduced. At the narrower bandwidth it is approximately $2.37$ USD.
These values expose approximation interacting with a weak denominator.
They are not expectations of an upstream ratio estimator and do not show
that an actual package returns these values.~\citep{exec17_s4_math}

Knowing the exact curvature in this artificial law would permit exact
population correction. In data, estimating curvature creates another random
quantity, and a weak first stage remains a sampling problem after correction.
The product should distinguish an approximation certificate, an interval for
the ratio, and a scientific assessment of exclusion and local compliance.
It should preserve unbounded or disconnected uncertainty where the selected
valid procedure permits it. A finite number from a stabilized denominator is
not itself evidence for a bounded causal effect.

\section{Finite contrasts, score support and portfolio transport}
\label{exec17:s4-support}

A second approximation arises when software evaluates predictions at nearby
scores. In the selected pinned CausalPy regression-discontinuity path, the
contrast uses design rows for treatment above $+\epsilon$ and control below
$-\epsilon$. This is a finite contrast, not automatically the limiting
cutoff effect. For example, let the two mean functions be
\[
 m_1(s)=18+5s+12s^2,\qquad m_0(s)=10+3s.
\]
Their cutoff difference is $8$ USD, while
\[
 m_1(\epsilon)-m_0(-\epsilon)=8+8\epsilon+12\epsilon^2.
\]
At $\epsilon=1/10,1/100,1/1000$, the contrasts are $8.92$, $8.0812$ and
$8.008012$ USD. Decreasing $\epsilon$ reduces this particular deterministic
approximation. It does not create data near the cutoff or eliminate model
misspecification. A numerical tolerance should therefore be expressed in
score units and translated into a justified monetary approximation bound
when the model permits one.~\citep{exec14_causalpy_primary,exec17_s4_math}

The same CausalPy source distinguishes expected mean predictions from
posterior predictive observations. Predictive variation can include outcome
noise; posterior uncertainty in a mean contrast is a different object.
The inspected scikit-learn route's singleton draw axes are not a posterior
sample. A common result schema should state the uncertainty object and
supported backend. Admitting an array shape as evidence of a distribution
would turn a formatting convention into an unsupported scientific claim.
\citep{exec14_causalpy_primary,exec17_s4_paths}

Discrete eligibility scores can remove the limiting comparison entirely.
Suppose $S$ takes only $-1$ and $1$, action is $1\{S\geq0\}$, and observed
outcomes are respectively $10$ and $14$ USD. Two response systems agree on
every observation:
\[
 \text{system I: }Y(a)=10+4a,\qquad
 \text{system II: }Y(a)=12+2S.
\]
Both have smooth extensions as functions of the score. The cutoff treatment
effect is $4$ in the first system and $0$ in the second. Continuity alone
cannot distinguish them because the score support never approaches zero.
Fitting a line across the gap chooses an extrapolation restriction. A
mass-point diagnostic or a bandwidth adjustment can reveal or respond to
the observed support, but cannot supply the unobserved response surface.
\citep{exec17_s4_math}

Possible remedies depend on the bank's decision. A defensible local
randomization model for a specified finite window would be a different
scientific assumption. A justified functional-form or curvature restriction
could support extrapolation or bounds. A redesigned randomized encouragement
could provide additional intervention information. Without such information
or restrictions, the correct product result is an unresolved effect set,
not a universal recommendation to use a different smoother.

Even a well-identified continuous-score cutoff effect need not equal a
portfolio effect. Suppose $\tau(0)=8$ USD and stipulate the sensitivity
restriction
\[
 |\tau(s)-\tau(0)|\leq L|s|,\qquad
 L=6\ \text{USD per score unit}.
\]
For a target population with equal mass at $s=-1/2$ and $s=1/2$, integrating
the pointwise envelope gives
\begin{equation}
 5\leq E_T\{\tau(S)\}\leq11\quad\text{USD}.
 \label{exec17:s4-transport-bounds}
\end{equation}
These bounds are sharp relative to the stipulated cutoff, sensitivity
restriction and observed sharp-threshold law. To construct witnesses, take
the observed arm mean to be $10+2s+8a$ with $a=1\{s\geq0\}$. Preserve
that mean on the observed arm, set the effect to $8-6|s|$ or $8+6|s|$,
and fill in the unobserved arm mean by subtracting or adding the effect.
The two continuous response systems agree on every observed mean and at the
cutoff while attaining the two portfolio endpoints.~\citep{exec17_s4_math}

At a cost of $7$ USD per action, the portfolio net-value set is $[-2,4]$ USD.
The local cutoff net value of $1$ USD does not determine its sign. The constant
$L$ is a scientific sensitivity parameter, not an estimate certified by
the local effect. Additional justified shape restrictions, intervention data
away from the cutoff or a narrower portfolio could reduce the set. A
posterior extrapolation can also be reported, provided its model-dependent
meaning remains explicit. General policy eligibility and transport cannot be
inferred from a local ratio alone.

\section{Adoption choices and a public comparison that can be built}
\label{exec17:s4-adoption}

The practical comparison is among scientific routes as well as languages.
DoubleML offers an inspected Python score architecture; \texttt{grf} supplies
R interfaces backed by native forest components; CausalML supplies heterogeneous
prediction and selected uncertainty paths; CausalPy supplies model-based
threshold predictions; and \texttt{rdrobust} supplies specialized local
polynomial procedures. An existing native component is evidence of an
implementation route, not evidence that a newly ported common core has the
same target or interval behavior. The following decisions are prospective
and scoped to the inspected paths.~\citep{exec17_s4_paths,exec17_s4_design}

\begin{longtable}{P{.19\textwidth}P{.35\textwidth}P{.36\textwidth}}
\toprule
Candidate route&Proposed use&Condition or refusal\\ \midrule
DoubleML IRM&Wrap selected ATE/ATTE score paths; retain callable-score extension
as a separate development.&Bind folds, nuisance predictions, clipping,
normalization and target. A new overlap score needs its own derivative,
remainder and statistical study.\\
\texttt{grf}&Wrap the supported average-effect and heterogeneous-prediction
routes before considering a native port.&Preserve target/method control flow,
cluster units and residual-regression semantics. Do not equate pointwise
forest intervals with selected-policy uncertainty.\\
CausalML&Adopt selected conditional predictions as candidate nuisances or
action scores.&Build a common-population evaluator for multiarm comparisons;
verify joint covariance. Retain the selected package variance separately when
its population differs.\\
CausalPy&Use the inspected threshold construction for explicit model-based
comparisons.&Name finite $\epsilon$, mean versus predictive uncertainty and
backend. Reject an unsupported cutoff interpretation for a score gap.\\
\texttt{rdrobust}&Compare specialized sharp/fuzzy local-polynomial procedures
on their admitted designs.&Bind the exact implementation, bandwidths,
polynomial orders and variance method; distinguish mass-point handling from
scientific identification.\\
Licensed Stata host&Offer a documented host route where the bank already has
the required entitlement.&The inspected package requires Stata 16 or later.
GPL-3 package rights and commercial host rights are separate; deployment,
redistribution and cost need a matched review.\\ \bottomrule
\end{longtable}

The inspected Stata \texttt{rdrobust} help declares version 11.1.1 dated
1 October 2026, with controls for outcome and bias-polynomial orders, effect
and bias bandwidths, fuzzy receipt, mass points and variance. The inherited
repository selection used a UTC end-of-day cutoff query. Whether that exact
revision was available by the strict Shanghai client cutoff remains unresolved.
The subsequently acquired mutable Stata 19 manuals also lack certified
cutoff-date bytes. Their inspected documented procedures can inform a bounded
comparison while those identity limits remain visible; they cannot establish
an exact as-of-date commercial capability census.~\citep{exec13_ivrd_rdrobust,exec13_ivrd_stata}

The commercial comparison should therefore ask whether the licensed route
answers the same target with the same support and sampling assumptions, which
diagnostics and uncertainty procedures it documents, and what integration and
operating costs apply. The proposal's desired advantage is preservation of
these distinctions across discovery, identification, estimation and reporting.
That is a proposed capability, not a measured victory over a commercial tool.
No proprietary price, actual deployment entitlement or new vendor runtime is
established by the present source review.

A bounded public/generated prototype can test the proposed advantage without
waiting for bank data. First freeze the package revisions, environments,
complete dependency rights and target contracts. Replay the exact weighted,
multiarm, resource and threshold laws above to test population arithmetic,
score conventions and expected refusals. Use the same declared target when
comparing algorithms; deliberately changed targets form separate cases.
The current exact reference calculations perform only the mathematical
part of that work.~\citep{exec17_s4_design}

For finite-sample statistical claims, separately register a prospective study:
sampling and assignment laws, outcome noise, independent units, nuisance
learners, folds, bandwidth rules, all uncertainty procedures, seed, fixed
request counts and complete failure accounting. Public covariates can supply
realistic margins while generated potential outcomes preserve known causal
truth. Their public origin does not make a generated intervention mechanism
authentic bank evidence. Compare common-target bias and error, interval
coverage and width, availability, allocation feasibility, policy loss and full
observed computation cost. Include clipped propensities, weak receipt jumps,
discrete support and dependent-unit mismatches as declared stress regimes,
rather than silently excluding failed requests.

Successful arithmetic and statistical comparisons would justify a bounded
prototype and identify where wrapping suffices or porting merits investment.
They would leave real confounding, score manipulability, feasible bank actions,
institutional capacity, acceptable loss and actual total prices to the later
qualification decision. Where two scientific models still agree on the
observations but disagree on the effect, more software cannot resolve the
ambiguity. The product must either acquire the missing information, narrow
its assumptions or preserve the effect set through the financial decision.

The next analysis develops the shared-control, allocation and discrete-score
cases further. The source comparisons then ask whether the implementations
preserve those targets and uncertainties. This order separates what the bank
needs to estimate from what a particular numerical interface happens to return.
\label{exec17:s4-end}


% SOURCE docs/release/2026-10-29/source/docs/writing-repair/2026-10-29/src/63bdca02c012ec17.tex CHAR 0
\section{Pair selection can reverse the joint uncertainty comparison}
\label{exec17:s4-common-population}

The shared-control calculation becomes more informative when the action groups
have different covariate composition. Consider two equally prevalent account
strata, $X=0,1$, with the following stipulated thirty-day law:
\[
\begin{array}{c|ccc|cc}
 X&e_0&e_1&e_2&\tau_1\ ({\rm USD})&\tau_2\ ({\rm USD})\\ \hline
 0&1/2&1/4&1/4&0&12\\
 1&1/4&1/4&1/2&20&0
\end{array}
\]
Let every arm have conditional residual variance $4$ USD$^2$, with
mean-zero residuals, and retain conditional exchangeability and support.
For example, independent residuals taking values $-2,2$ with equal
probability realize that variance. Baseline means may be fixed at any
declared common values. The effects averaged over all eligible accounts are
$\theta_1=10$ and $\theta_2=6$ USD. This complete observable-law specification
permits exact score calculations; it is not a fitted bank response model.
\citep{exec17_s4_extension}

Restricting observations to a treatment/control pair changes the covariate
law. Write $B_j=1\{A\in\{0,j\}\}$ and
$s_j(x)=e_0(x)+e_j(x)$. Bayes' rule gives
\begin{equation}
 P(X=x\mid B_j=1)=\frac{P(X=x)s_j(x)}{P(B_j=1)},\qquad
 \theta_j^{\rm pair}=
       \frac{E\{s_j(X)\tau_j(X)\}}{E\{s_j(X)\}}.
 \label{exec17:s4-pair-target}
\end{equation}
For action 1, the retained mass is $5/8$ and the high-stratum share is
$2/5$, giving $\theta_1^{\rm pair}=8$ USD. For action 2, retained mass is
$3/4$ and both strata remain equally prevalent, giving
$\theta_2^{\rm pair}=6$ USD. The pair contrast for action 1 and the
full-population contrast differ even with the outcome regressions known
exactly. Dropping the other action is therefore not an innocuous way to
define the evaluation population.

This calculation explains a specific concern in the selected CausalML
implementation. Its point mean averages predictions over the full input
matrix, while its analytic uncertainty calculation uses a pair mask and
pair-row denominator. The derived $8$ is an explanatory pair-population
mean, not a claim that the package returns a point estimate of $8$; with
the true predictions and the full balanced input, its described point
average is $10$. Nor does the selected variance formula automatically
become the valid variance for either estimand under arbitrary assignment.
The full-input mean and masked ingredients require a matched statistical
argument.~\citep{exec14_causalml_primary}

For the common-population scores in Equation~\eqref{exec17:s4-shared-control},
the heterogeneity variances are $100$ and $36$, and their covariance is
$-60$ USD$^2$. The expected shared-control noise contribution is
$E(4/e_0)=12$; the active-arm contributions are $16$ and $12$.
Consequently
\begin{equation}
 \Sigma=\begin{pmatrix}128&-48\\-48&60\end{pmatrix},\qquad
 \operatorname{Var}(\phi_2-\phi_1)=284\quad {\rm USD}^2.
 \label{exec17:s4-negative-joint}
\end{equation}
Here treating the two estimates as independent gives $188$, understating
the variance of their difference. This reverses the direction of the error
in the constant-effect example, where independence gave $48$ instead of
$32$. A common control creates positive residual covariance, but opposing
heterogeneous effects can dominate it.

There is a useful separate check. The conditional difference
$\tau_2-\tau_1$ takes values $12,-20$ and has variance $256$.
In $\phi_2-\phi_1$ the control residual cancels, leaving active-arm
contributions $12+16=28$. Their sum is $284$, agreeing with the matrix
calculation. An enumeration of the two strata, three observed arms and two
residual values gives the same covariance directly from twelve observable
atoms. Thus the common-population comparison can be verified without
treating a fitted prediction matrix as a covariance model.
\citep{exec17_s4_extension}

One possible reuse design keeps the candidate package's conditional
predictions while constructing a separate common-population score evaluator.
It must retain every eligible account's covariates, original multiarm
propensities, action coding and independent sampling unit. For an externally
specified evaluation population, a declared transport or reference-sample
argument is additionally needed. Merely multiplying a pair-specific standard
error by a population-size ratio does not restore the missing heterogeneity
or joint covariance. The proposed evaluator is a development route, whose
statistical and computational behavior still needs prospective comparison.

\section{The experiment design depends on which comparison matters}
\label{exec17:s4-design-criterion}

The covariance argument also changes how a future public or prototype
experiment should be designed. Suppose a new randomized study can choose
constant probabilities $p_0,p_1,p_2>0$, summing to one. Hold the target
population fixed, take constant effects and residual variance $4$ in each
arm, and consider iid observations. If the criterion is the sum of the
two ATE influence variances, it is
\begin{equation}
 T(p)=\frac{8}{p_0}+\frac{4}{p_1}+\frac{4}{p_2}.
 \label{exec17:s4-design-trace}
\end{equation}
The control variance is charged twice because it enters both comparisons.
For a sample of size $n$, divide this criterion by $n$. It measures the
chosen precision objective, not financial utility or an experiment's total
cost.

For positive coefficients $a_j$, Cauchy--Schwarz yields
\[
 \left(\sum_j\frac{a_j}{p_j}\right)\left(\sum_jp_j\right)
       \geq\left(\sum_j\sqrt{a_j}\right)^2.
\]
Equality requires $p_j$ proportional to $\sqrt{a_j}$. Applying this with
$a=(8,4,4)$ gives
\[
 p_0=\sqrt2-1,\qquad p_1=p_2=(2-\sqrt2)/2,\qquad
 T_{\min}=24+16\sqrt2\simeq46.63\quad {\rm USD}^2.
\]
The earlier design $(1/2,1/4,1/4)$ gives $48$. Symmetry provides another
derivation: set $p_1=p_2=(1-p_0)/2$ and minimize
$8/p_0+16/(1-p_0)$. Its second derivative is
$16/p_0^3+32/(1-p_0)^3>0$, and its stationary point is $\sqrt2-1$.
Both arguments certify the global optimum for this stated criterion.
\citep{exec17_s4_extension}

If the decision is instead which active action is better, the relevant
constant-effect influence variance is
\[
 D(p)=4/p_1+4/p_2.
\]
It contains no control residual. Minimizing it alone would drive the
control probability toward zero, sacrificing precision and ultimately
support for the two comparisons against no action. If the study separately
requires $p_0\geq1/5$, this criterion is minimized at
$(p_0,p_1,p_2)=(1/5,2/5,2/5)$, giving $D=20$ but $T=60$.
The earlier design gives $D=32$ and $T=48$. Neither design dominates the
other on both precision objectives.

The distinction is useful for a bank deciding whether to fund a pilot.
An experiment intended to choose between two already acceptable servicing
actions has a different information objective from one intended to show
that either beats no contact. If effects are heterogeneous, their fixed
population covariance terms add to these formulas; if variances differ
by stratum, costs differ by arm, or randomization is clustered, the design
problem must be reformulated. The displayed optimum assumes equal per-row
resource cost and the specific independent law. It does not authorize a
real randomized bank intervention or assert power for an unspecified loss
threshold.

\section{Effect sets lead to different robust allocation objectives}
\label{exec17:s4-robust-allocation}

An effect need not be reduced to one number before allocation. Let there
be finitely many strata with known masses $w_i$, a fixed feasible allocation
set $\mathcal D$, and a rectangular net-gain class
$\mathcal G=\prod_i[\ell_i,u_i]$. This class is stipulated scientific
sensitivity information. A statistical confidence set would need a separate
simultaneous coverage argument, and correlations or common structural
restrictions could make a rectangle unnecessarily broad. For fixed known
costs, an effect interval can be translated to a net-gain interval by
subtracting the cost from both endpoints.

Write $V_g(d)=\sum_iw_i g_i d_i$, where $d_i\geq0$. The maximin policy
maximizes its worst possible value. Rectangularity and nonnegative
allocations imply
\begin{equation}
 \inf_{g\in\mathcal G}V_g(d)=\sum_iw_i\ell_i d_i.
 \label{exec17:s4-maximin}
\end{equation}
Thus maximin allocation can use the earlier resource optimizer with lower
gains. Its objective is a guarantee against the declared class, not an
estimate that every gain equals its lower endpoint.

Minimax regret asks a different question: how much would the chosen policy
lose relative to the best feasible policy if the true gain vector became
known? For a fixed $d$, its worst regret is
\begin{align}
 R(d)&=\sup_{g\in\mathcal G}\sup_{z\in\mathcal D}
                   \sum_iw_i g_i(z_i-d_i)\notag\\
 &=\sup_{z\in\mathcal D}\sum_iw_i\left\{
       u_i(z_i-d_i)_+ +\ell_i\min(z_i-d_i,0)\right\}.
 \label{exec17:s4-robust-regret}
\end{align}
The two suprema commute because they range over a product set. For each
fixed difference, the gain maximizing its contribution is the upper endpoint
for a positive difference and the lower endpoint for a negative one. The
oracle and chosen policy must share the same feasible class. Allowing the
oracle extra capacity would measure a different loss.

Take two equally prevalent strata, one resource unit per action, budget
$B=1/2$, and net gains
\[
 g_1\in[2,8],\qquad g_2\in[4,6]\quad {\rm USD}.
\]
All gains are positive, so it suffices to allocate the full budget. Write
$d=(t,1-t)$, $0\leq t\leq1$. Maximin value is
$\tfrac12\{2t+4(1-t)\}=2-t$, maximized at $t=0$ with guarantee $2$ USD.
If stratum 1 is better, regret is
$\tfrac12(1-t)(g_1-g_2)$; if stratum 2 is better, it is
$\tfrac12t(g_2-g_1)$. Therefore
\begin{equation}
 R(t)=2\max(1-t,t),\qquad t_{\rm minimax}=1/2,\quad R(1/2)=1\ {\rm USD}.
 \label{exec17:s4-robust-example}
\end{equation}
The endpoints $(g_1,g_2)=(8,4)$ and $(2,6)$ attain the two regret branches.
Their equality at $t=1/2$ and the bound
$\max(t,1-t)\geq1/2$ prove optimality. The mixed policy's guaranteed value
is $3/2$ USD. The maximin policy guarantees more, $2$ USD, while having
worst regret $2$ USD. Different objectives produce different defensible
recommendations from the same causal information.~\citep{exec17_s4_extension}

The interval midpoints both equal $5$ USD. A point-value optimizer can
therefore break that tie arbitrarily and choose either pure allocation,
each with worst regret $2$. The minimax mixture halves that worst regret.
This does not show that midpoint estimation is universally inferior:
a justified probability model for gains would support a distinct expected
utility criterion. It shows why the proposal should expose the decision
objective and preserve causal ambiguity instead of silently turning an
effect set into its midpoint. The mixture concerns population fractions;
an actual finite roster still requires an admissible implementation of
capacity and randomization.

\section{A proved finite-policy evaluation bound}
\label{exec17:s4-finite-policy}

The deterministic regret comparison can be supplemented by a transparent
statistical result in a restricted experiment. Suppose a training sample
produces a fixed finite collection of $M$ policies $d_1,\ldots,d_M$ and
outcome predictions $\widetilde m_a$. Conditional on that training sample,
evaluate the policies on $n$ new iid accounts from the declared population.
Let treatment be randomized with known $e(X)$ bounded away from zero and
one, and let costs $c(X)$ be known. Define each account's net-value score
\begin{align}
 W_j={}&d_j(X)\left\{\widetilde m_1-\widetilde m_0-c(X)
       +\frac{A}{e(X)}(Y-\widetilde m_1)
       -\frac{1-A}{1-e(X)}(Y-\widetilde m_0)\right\}.
 \label{exec17:s4-policy-score}
\end{align}
Conditional on $X$ and training, the two residual terms have expectations
$m_1-\widetilde m_1$ and $-(m_0-\widetilde m_0)$. Thus
$E(W_j\mid\text{training})=V_j=E\{d_j(X)(\tau(X)-c(X))\}$.
No correctness of the trained outcome regressions is needed for this
expectation because the assignment probability is known and correct.
This argument does not carry over unchanged to a fitted observational
propensity or an invalid scientific design.

Assume every $W_j$ lies in a known interval $[a,b]$, of width $R=b-a$.
The following elementary exponential bound gives a fully stated route to
simultaneous evaluation. For any $Z\in[0,1]$ with mean $p$, convexity of
the exponential gives
\[
 E e^{t(Z-p)}\leq (1-p)e^{-tp}+p e^{t(1-p)}.
\]
The logarithm of the right side has value and derivative zero at $t=0$;
its second derivative is the variance of a tilted Bernoulli variable,
at most $1/4$. Integrating that bound twice gives a logarithm at most
$t^2/8$. Scaling $Z=(W_j-a)/R$, independence of the evaluation accounts
therefore implies
\[
 E\exp\{\lambda\textstyle\sum_{i=1}^n(W_{ij}-V_j)\}
       \leq\exp(n\lambda^2R^2/8).
\]
The exponential Markov inequality followed by minimization at
$\lambda=4u/R^2$ yields
$P(\overline W_j-V_j\geq u)\leq\exp(-2nu^2/R^2)$.
Applying it to the negative deviation and taking a union bound over the
$M$ policies proves, conditionally on training,
\begin{equation}
 P\left(\max_j|\overline W_j-V_j|\leq
       R\sqrt{\frac{\log(2M/\alpha)}{2n}}\right)\geq1-\alpha.
 \label{exec17:s4-uniform-bound}
\end{equation}
The scores for different policies may be correlated on the same evaluation
accounts; this union argument does not require independence across policies.
It requires independence across evaluation accounts and the fixed conditional
candidate set.~\citep{exec17_s4_extension}

Let $\widehat j$ maximize evaluated mean value over that set, and let
$j^*$ maximize true value. On the simultaneous event, writing its radius
as $b_n$, one obtains
\[
 V_{j^*}-V_{\widehat j}
 \leq (V_{j^*}-\overline W_{j^*})
      +(\overline W_{\widehat j}-V_{\widehat j})\leq2b_n.
\]
An optimization tolerance $\eta$ adds $\eta$. Selection using this evaluation
set is allowed because the event holds for every member of the finite,
conditionally fixed collection. Continually inventing new policies from
these evaluation outcomes, or applying the result to an unrestricted
adaptive learner class, would exceed the proof's domain.

For an explicit monetary scale, take $e=1/2$, outcomes and trained arm
means in $[0,30]$ USD, costs in $[0,5]$ USD and $0\leq d_j\leq1$.
When $A=1$, the bracket in Equation~\eqref{exec17:s4-policy-score} is
$2Y-\widetilde m_1-\widetilde m_0-c$; when $A=0$, it is
$\widetilde m_1+\widetilde m_0-2Y-c$. Both lie in $[-65,60]$ USD;
multiplication by a policy fraction preserves that interval. Thus $R=125$.
For $M=16$, $\alpha=1/20$ and $n=10{,}000$, the simultaneous radius is
approximately $2.2468$ USD and the selected-policy regret guarantee is
$4.4936$ USD. These are proved upper bounds under the stipulated design,
not observed interval widths or measured regret.

To make the same sufficient regret guarantee at most $1$ USD, without
optimization error, requires
\begin{equation}
 n\geq 2(125)^2\log(640),\qquad n_{\rm integer}=201{,}921.
 \label{exec17:s4-evaluation-size}
\end{equation}
The conservative requirement is not a claim that this is the necessary
sample size or that a bank should purchase so many observations. A sharper
variance-sensitive analysis, fewer candidates, better score bounds or a
different acceptable loss may change the design substantially. Boundedness
is substantive: unbounded monetary outcomes, unknown truncation bias,
dependent accounts, estimated assignment probabilities or altered policy
classes need their own justified analysis. An empirical coverage study can
compare suitable procedures, but cannot be replaced by counting the exact
reference cases above.

For a product proposal, this calculation explains why public covariates and
generated interventions are useful before a costly authentic pilot. They can
exercise score conventions, evaluation splits, policy selection and complete
failure accounting while the causal truth is known by construction. The
costly study's merit should be judged against a declared precision or loss
goal. Neither the bound nor a generated study establishes the actual bank's
assignment law, feasible action set or acceptable monetary loss.

\section{What scientific information can close a discrete-score gap?}
\label{exec17:s4-gap-bounds}

The earlier two-point example shows that smoothness without a quantitative
restriction does not identify the cutoff effect. A more informative answer
can be obtained if the sponsor can defend explicit slope bounds. Suppose
the only observed scores are $-a$ and $b$, where $a,b>0$, sharp assignment
holds, and the observed means are $y_-$ and $y_+$. Let
$\mu_j(s)=E\{Y(j)\mid S=s\}$ be potential-mean extensions to the cutoff.
Stipulate that $\mu_0$ and $\mu_1$ are respectively Lipschitz with known
constants $L_0,L_1$ on the relevant intervals. Consistency then implies
\[
 \mu_0(0)\in[y_- -L_0a,y_-+L_0a],\qquad
 \mu_1(0)\in[y_+-L_1b,y_++L_1b].
\]
Subtracting the intervals gives
\begin{equation}
 \tau(0)\in[y_+-y_- -L_1b-L_0a,
             y_+-y_- +L_1b+L_0a].
 \label{exec17:s4-gap-sharp}
\end{equation}
These are sharp bounds for this two-point observed law and these separate
Lipschitz restrictions, absent further outcome or shape constraints.
For the lower endpoint choose the linear extensions
\[
 \mu_1(s)=y_++L_1(s-b),\qquad
 \mu_0(s)=y_-+L_0(s+a).
\]
They preserve the observed arm means, obey the slope bounds and attain the
lower difference at zero. Reversing both slopes attains the upper endpoint.
Convex mixtures preserve the same observations and restrictions while
attaining every intermediate difference. This is a constructive scientific
identified set, not a confidence interval produced by a smoother.
\citep{exec17_s4_extension}

For $a=b=1$, $y_-=10$, $y_+=14$ and $L_0=L_1=1$, the effect set is
$[2,6]$ USD. With action cost $3$ USD the net set is $[-1,3]$. A unit
slope restriction excludes the earlier zero-effect response system, whose
baseline slope was two; that exclusion comes from the additional scientific
restriction. At $L_0=L_1=2$, the set widens to $[0,8]$ and includes both
earlier systems. The data have not changed. The quantitative assumption
explains exactly why the conclusion changes.

More observations at the same two scores may reduce sampling error in
$y_-$ and $y_+$, but cannot remove the deterministic width
$2(L_1b+L_0a)$. Closer supported scores, justified smaller slope bounds,
intervention information in the gap or an alternative scientifically
defensible target could reduce that ambiguity. When observed means are
estimated, a valid joint confidence region for them can be propagated
through the endpoint expressions; its sampling uncertainty remains distinct
from the stipulated $L_j$. A Bayesian extrapolation can likewise be useful
with an explicit prior and model, while answering a different uncertainty
question.

The effects proposal can consequently offer concrete remedies with precise
limits: align evaluation populations, preserve joint comparisons, optimize
the intended resource and uncertainty criterion, and state which additional
information narrows an unresolved effect set. Selected software can support
parts of those routes. Their integration, runtime qualification and bank
assumptions remain explicit development and adoption obligations.
\label{exec17:s4-extension-end}


% SOURCE docs/release/2026-10-29/source/docs/writing-repair/2026-10-29/src/a86069a62dfe2036.tex CHAR 0
\section{R effect estimation: preserve the target before choosing the learner}
\label{exec20:rfx-scope}

A servicing team may ask which customers should receive an intervention, what
the intervention achieves on average, or how much additional credit should be
offered. These questions require different statistical objects. Let $X$ denote
information available before the decision, $A\in\{0,1\}$ a recorded intervention,
and $Y(a)$ its potential monetary outcome. The conditional effect
$\tau(x)=E\{Y(1)-Y(0)\mid X=x\}$, its population average, and its average among
customers actually treated are not interchangeable. For a continuous amount $W$,
the dose-response curve $\mu(w)=E\{Y(w)\}$ introduces another target. Identification
still requires the relevant consistency, exchangeability and support assumptions;
orthogonal nuisance estimation does not supply them. Forest and double-machine-
learning results impose their own sampling and regularity conditions.
\citep{dml2018,grf2018}

The R routes examined here have distinct strengths. The selected \texttt{grf}
interface exposes conditional forests and several average-effect targets;
\texttt{tmle} addresses observed binary point treatments and outcome missingness;
R \texttt{DoubleML} exposes explicit orthogonal-score models. \texttt{CausalGPS}
and \texttt{causaldrf} supply different paths to continuous-exposure analysis.
The proposed product should wrap their qualified statistical contracts before
considering a native port. Matching an algorithm name or programming language
does not establish matching targets or uncertainty.
\citep{exec13_grf_rpath,exec20_rfx_tmle,exec20_rfx_doubleml_r,
exec20_rfx_causalgps,exec20_rfx_causaldrf}

\subsection{An R score is a specification, not a Python alias}

The inspected R DoubleML revision is dated 1 September 2026 and declares
development version \texttt{1.0.1.9000}. Its R6 and \texttt{mlr3} implementation
must be assessed separately from the Python installation considered earlier.
The selected interactive-regression class requires one binary treatment. Write
$m(x)=P(A=1\mid X=x)$ and $g_a(x)=E(Y\mid A=a,X=x)$. Its ATE score uses
\begin{equation}
 \psi_b=g_1(X)-g_0(X)+\frac{A\{Y-g_1(X)\}}{m(X)}
 -\frac{(1-A)\{Y-g_0(X)\}}{1-m(X)},\qquad \psi_a=-1,
 \label{exec20:rfx-irm}
\end{equation}
and solves the empirical equation $\overline{\psi_a}\theta+
\overline{\psi_b}=0$. For ATTE, the selected code instead uses the control
outcome nuisance, a treated fraction calculated within each test fold, and
$\psi_a=-A/\widehat p$. It fits the treated-outcome nuisance for ATE and custom
scores, but does not require that fit for its built-in ATTE path. These choices
must survive the wrapper: changing the target cannot be implemented merely by
changing a returned field name. The propensity truncation option clips predictions;
it does not itself define a new trimmed population.
\citep{exec20_rfx_doubleml_r,dml2018}

The nuisance helper supplies explicit training and test indices to custom
\texttt{mlr3} resampling. Treatment-specific outcome fits restrict training
indices while retaining the requested prediction rows. This is distinct from
cross-validation used to tune a learner: the inspected tuning method defaults
to global tuning and has an explicit option for tuning separately within folds.
A proposed validation protocol must record both sets of folds and any use of
outcomes during tuning. Neither a default option nor the existence of custom
scores proves that an arbitrary nuisance learner or new score satisfies a theorem.
\citep{exec20_rfx_doubleml_r,dml2018}

For continuous $W$, the partialling-out score solves a different moment. Let
$e(X)=E(W\mid X)$, $\ell(X)=E(Y\mid X)$ and $v=W-e(X)$. Its population target,
when the denominator is positive, is
\begin{equation}
 \theta_{\mathrm{PLR}}=
 \frac{E[v\{Y-\ell(X)\}]}{E(v^2)}.
 \label{exec20:rfx-plr}
\end{equation}
Under the stipulated conditional model
$Y=a(X)+\beta(X)W+\epsilon$, with $E(\epsilon\mid W,X)=0$, conditioning shows
$E(vY\mid X)=\beta(X)\operatorname{Var}(W\mid X)$. Hence
\begin{equation}
 \theta_{\mathrm{PLR}}=
 \frac{E\{\operatorname{Var}(W\mid X)\beta(X)\}}
 {E\{\operatorname{Var}(W\mid X)\}},
 \label{exec20:rfx-weighted}
\end{equation}
rather than generally $E\{\beta(X)\}$. Two equally probable risk groups with
slopes 1 and 3 and treatment variances 1 and 4 give respectively $13/5$ and 2.
The difference is a population-target difference even with perfect nuisances.
These are exact authored laws, not measurements of package error.
\citep{exec20_rfx_doubleml_r,exec20_rfx_reference}

\subsection{Forests, clustering and the information needed for uncertainty}

In the selected \texttt{grf} R path, all-sample AIPW averaging uses scores and
sums weighted centered scores within declared clusters for its variance.
The continuous-treatment score learns conditional treatment variance with an
additional forest and divides a treatment residual by that variance. Positivity
of this denominator and stable estimation are substantive conditions. A returned
local slope cannot be relabelled as an arbitrary finite-dose effect. The overlap
target uses a separate residual-regression path; the documented method choice
does not bypass earlier guards. In particular, the inspected average-effect
function rejects TMLE with clusters or sample weights. Honesty in forest
construction, treatment residualization and independent-cluster inference solve
different problems.
\citep{exec13_grf_rpath,grf2018}

The selected R DoubleML base class draws partitions over cluster identifiers
and has separate one- and two-way cluster variance calculations. This does not
establish validity for arbitrary serial dependence or few calendar periods.
Moreover, its multiplier-bootstrap method explicitly stops for clustered data.
Its joint-interval method requires bootstrap output, whereas the ordinary
interval method uses a normal critical value with the reported standard error.
A wrapper must preserve this unsupported combination; silently substituting
row multipliers would change the sampling law. A cluster-aware pointwise
standard error and a jointly valid set of comparisons are different deliverables.
\citep{exec20_rfx_doubleml_r}

For servicing data, the request must therefore name the independent sampling
units, repeated-customer identifiers, common calendar shocks, and how entire
units are kept out of nuisance-training folds. The public prototype can compare
supported regimes with known generated truth before an expensive bank trial.
It must retain failed folds, absent treatment arms, unsupported uncertainty
requests and denominator failures. These are proposed admission requirements;
the source inspection does not execute either R engine.

\subsection{Missing outcomes do not make missing treatment labels observed}

The pinned \texttt{tmle} 2.1.1 help text describes outcome missingness, but its
\texttt{Delta} argument also mentions missing treatment assignment. The selected
function body provides a narrower operational contract: its argument verifier
requires no missing values in $A$ or $W$, and the public entrypoint stops when
verification fails. Earlier preprocessing can also fail on an incompletely
observed treatment vector. Passing \texttt{A=NULL}, or an all-zero vector,
selects a constant-treatment path by replacing $A$ with ones. The documentation
describes the no-variation case as mean-outcome estimation. This is not treatment
imputation, and the missingness regression still conditions on recorded $A$.
The proposal should reject arbitrary missing-$A$ admission for this inspected
path and retain the conflicting wording for clarification and native fixtures.
\citep{exec20_rfx_tmle}

There is a scientific reason to keep those fields separate. Consider an authored
law in which observation group $R$ and randomized treatment $A$ are independent
fair coins. When $R=1$, both worlds have $Y(a)=a$ and the treatment is recorded.
When $R=0$, the treatment label is erased: world I has $Y(a)=a$, while world II
has $Y(a)=1-a$. The observed tuple consists of $R$, the treatment label when
present, and $Y$. Both worlds give the same distribution of that tuple: each
observed group has a fair binary outcome. Nevertheless,
\begin{equation}
 E\{Y(1)-Y(0)\}_{\mathrm{I}}=1,
 \qquad E\{Y(1)-Y(0)\}_{\mathrm{II}}=0.
 \label{exec20:rfx-missing}
\end{equation}
This witness deliberately lacks treatment-observation positivity in group
$R=0$ and grants no restriction transporting effects from $R=1$. It does not
claim impossibility under every missingness model. At a stipulated USD 100
outcome scale, the two effects are USD 100 and USD 0 despite the identical
observations. More rows with the same missing information cannot distinguish
them; verified labels or justified additional restrictions can.
\citep{exec20_rfx_reference}

\subsection{Dose curves require units, support and complete transformations}

A conditional regression slope also differs from a dose curve in one homogeneous
group. Suppose incremental credit $W$ is measured in thousands of dollars,
is uniform on $[0,2]$, and the stipulated margin curve is
$Y(w)=20w^2-5w$ dollars. Direct polynomial integration gives
\begin{equation}
 \frac{\operatorname{Cov}(W,Y)}{\operatorname{Var}(W)}=35,
 \quad Y'(0)=-5,\quad Y(1)-Y(0)=15.
 \label{exec20:rfx-dose}
\end{equation}
Using the projection slope to value a USD 1,000 increase from zero would report
USD 35 in this law, whereas the stipulated contrast is USD 15. Expressing credit
in dollars changes the slope to $35/1000$; it does not change the finite monetary
contrast. This example illustrates target and unit errors, not bank profitability.
\citep{exec20_rfx_reference}

The inspected \texttt{CausalGPS} kernel path standardizes an exposure residual
and stores an estimated residual density. Its weighting helper divides an
estimated marginal exposure density by that stored quantity. Under a
location-scale model $W=e(X)+s(X)U$, with $s(X)>0$, the change of variables is
\begin{equation}
 f_{W\mid X}(w\mid x)=
 \frac{1}{s(x)}f_{U\mid X}\!\left(\frac{w-e(x)}{s(x)}\middle|x\right).
 \label{exec20:rfx-jacobian}
\end{equation}
The inspected kernel branch does not divide by $s(x)$ before returning its GPS.
With different conditional scales, replacing the original-unit density with
residual density changes relative weights. An exact two-row illustration with
scales 1 and 2, the same residual density and marginal numerator, gives weights
1 and 2 rather than 1 and 1. Applied to stipulated outcomes USD 0 and USD 100,
their weighted means are USD $200/3$ and USD 50. This is algebraic motivation
for a conformance fixture, not a measured error for every package output.
\citep{exec20_rfx_causalgps,exec20_rfx_reference}

Under a constant positive change of exposure units, both a correctly expressed
marginal density and conditional density acquire the same reciprocal unit
factor, so their ratio is invariant. A covariate-dependent omitted scale is a
different issue. A public test should vary units and conditional scales separately,
then trace the resulting weighting and response fits. The selected response
routine supplies parametric, semiparametric and nonparametric routes; it rejects
negative weights but changes an all-zero weight vector to equal weights.
The product should refuse an all-zero causal weighting request before that
fallback changes its meaning. Balance statistics and effective sample size are
useful diagnostics, not proof of exchangeability or calibrated uncertainty.
Trimming and complete-case filtering must also be reflected in the target
population and reported support.
\citep{exec20_rfx_causalgps}

The selected \texttt{causaldrf} \texttt{hi\_est} path follows a different design.
It fits a treatment model, fits an outcome regression on treatment and the
estimated score, then recalculates the score at each requested dose for every
row before averaging predicted outcomes. This counterfactual-score step must
be preserved. Its log and square-root branches evaluate densities on transformed
scales without applying the corresponding original-scale Jacobian in that
function. At a fixed positive dose those Jacobian factors are constant across
covariates, so their omission does not by itself prove a failure of a balancing
argument at that dose. It does change the density interpretation and the
representation supplied to a restricted response regression across doses.
Positive support, transformation rules and the whole fitted response surface
therefore need explicit comparison before adoption.
\citep{exec20_rfx_causaldrf}

The accompanying \texttt{get\_ci} helper receives a matrix of bootstrap
coefficients, forms linear or quadratic predictions, and sorts each dose column.
It neither generates the resamples nor refits the treatment model. Its separate
column order statistics are pointwise constructions, not a simultaneous curve
band. The selected lower index is $\lfloor0.025B\rfloor$, which is zero when
$B<40$; native edge cases need inspection rather than a claim that increasing
$B$ solves statistical validity. A proposed curve experiment must specify which
sampling units are resampled, which stages are refitted, and whether it evaluates
pointwise coverage, joint coverage, or a selected optimal-dose decision.
\citep{exec20_rfx_causaldrf,exec20_rfx_reference}

\subsection{A bounded reuse proposal and the remaining comparison}

For a first public prototype, the candidates are R DoubleML's supported
IRM/PLR score paths, \texttt{grf}'s target-specific forest paths, and
\texttt{tmle}'s observed point-treatment paths. Each remains conditional on
its complete statistical and software contract. The developed \texttt{grf}
review in Section~\ref{exec21:grfr-scope} follows nuisance provenance and
target-specific control flow. Section~\ref{sec:exec24-tmle-reuse} additionally
identifies a specific intermediate-variable risk-ratio bootstrap interval
mismatch, resample-count and missing-output conventions, and an unresolved
rights question. The affected interval is refused pending a reviewed correction;
these findings cannot be waived by the earlier point-treatment comparison.
Requests should carry target, units, treatment-label status, support, weighting,
folds and sampling units. Responses should preserve unsupported modes and
diagnostic failures. Native ports remain conditional on demonstrated semantic
and numerical parity; translating a wrapper does not reproduce a learner's
statistical contract.
\texttt{CausalGPS} remains a conditional candidate pending scale and downstream
weight checks, while \texttt{causaldrf} is a useful proposed parametric baseline
subject to transformation and full-pipeline uncertainty qualification.

The selected originals establish these specific contrasts, not a complete R
market census. Further candidates, including \texttt{bartCause},
\texttt{causalweight}, other continuous-exposure methods, and their aliases,
still require matched primary-source review. The Bayesian forest comparison
elsewhere in this chapter does not close these package-specific omissions. Package license declarations and
copyright notices do not settle all distribution or learner-dependency rights.
No R package has been executed in this source comparison. Public and generated
experiments can decide which wrappers merit a prototype; bank truth, deployment
acceptance and observed returns remain subsequent questions. The immediate
comparison must preserve the question being estimated before ranking a learner.


% SOURCE docs/release/2026-10-29/source/docs/writing-repair/2026-10-29/src/294eb15eb39615af.tex CHAR 0
\section{EconML reuse: preserve the target, history and sampling experiment}
\label{exec19:econml-scope}

A bank choosing which customers should receive a servicing offer needs an
effect for a specified population and decision time. A flexible treatment-effect
package can supply useful estimation machinery, but the meaning of the output
depends on the effect model, the variables used for adjustment and the population
over which predictions are averaged. A repeated-contact application adds a
second obligation: the computational history must preserve the customer's
actual sequence of decisions. These obligations motivate a selective EconML
reference profile, with a separate contract for a future executable adapter.

The primary source is the Python repository \texttt{py-why/EconML}.
The inspected commit is
\texttt{d625dd696ea28c62e6e4}\allowbreak\texttt{d8baea6e4154b98c30b4},
dated 1 October 2026. The source declares version 0.17.0, and its README records that version's
31 July 2026 release. Seventeen selected original files cover package notices,
dependency declarations, static DML, DynamicDML, common inference machinery
and selected upstream tests. The immutable source revision is the identity
used here; it does not authenticate every wheel or installation labelled
0.17.0. The selected files support source-level conclusions, while an actual
native fit and complete estimator-family qualification remain separate work.
\citep{exec19_econml_source}

\subsection{A residual projection can answer a different banking question}
\label{exec19:econml-target}

Write $A$ for an offer indicator, $Y$ for its monetary outcome, $X$ for
pre-action characteristics used to describe effect heterogeneity, and $W$ for
additional adjustment variables. In the selected DML formulation, the nuisance
functions are $m(X,W)=E[A\mid X,W]$ and $q(X,W)=E[Y\mid X,W]$.
Residuals $\widetilde A=A-m(X,W)$ and
$\widetilde Y=Y-q(X,W)$ enter the objective
\begin{equation}
 \widehat\theta\in\mathop{\rm arg\,min}_{\theta\in\Theta}
 \frac{1}{n}\sum_{i=1}^n
 \{\widetilde Y_i-\theta(X_i)\widetilde A_i\}^{2}.
 \label{exec19:econml-objective}
\end{equation}
The selected documentation states a structural residual relationship and
orthogonality conditions; the R-learner implementation forms the corresponding
residuals and final regression. A causal reading additionally needs the stated
identifying information and an appropriate effect-function class $\Theta$.
An accurate nuisance model alone cannot correct a deliberately inadequate
effect model. Variables admitted to $W$ also need causal justification:
automatic adjustment for a post-offer variable can change the question.
\citep{exec19_econml_source}

For a binary offer under a sufficient observed adjustment variable $S$, let
$e(S)=P(A=1\mid S)$ and $\tau(S)=E[Y(1)-Y(0)\mid S]$. With exact
nuisances and a constant final effect, differentiating the population squared
loss gives
\begin{equation}
 \theta_*=
 \frac{E[\widetilde A\widetilde Y]}{E[\widetilde A^2]}
 =\frac{E[e(S)\{1-e(S)\}\tau(S)]}
        {E[e(S)\{1-e(S)\}]}.
 \label{exec19:econml-projection}
\end{equation}
The second equality uses the identified binary-treatment conditional mean
relationship and a positive residual variance. Thus the constant projection
weights strata by treatment variation. The population average effect instead
uses the target population's stratum probabilities. They agree under a
correct constant-effect model, or particular balancing conditions, but need
not agree when effect heterogeneity has been omitted.
\citep{exec19_econml_reference}

Consider a fully specified example with $S\sim\mathrm{Bernoulli}(1/2)$,
$e(0)=1/2$, $e(1)=1/10$, and potential outcomes $Y(a)=100Sa$ in USD.
The offer is assigned using only observed $S$, so there is no hidden
confounding in this construction. The two strata have effects USD 0 and
USD 100. Direct enumeration gives
\begin{equation}
 E[Y(1)-Y(0)]=50,\qquad
 E[\widetilde A\widetilde Y]=\frac92,\qquad
 E[\widetilde A^2]=\frac{17}{100},\qquad
 \theta_*=\frac{450}{17}\simeq26.47.
 \label{exec19:econml-money}
\end{equation}
At an assumed offer cost of USD 35, offering to the whole target population
has expected net value USD 15 per customer. Labelling the constant projection
as that population effect would instead report net value
$450/17-35=-145/17$, about USD $-8.53$. This is a consequence of the
specified misspecification and target mismatch, not a demonstrated EconML
error under a correctly specified effect model. Including the effect modifier
in an adequate $\Theta$ and then averaging over the declared target population
is one route to the intended question in this example.
\citep{exec19_econml_reference}

The base \texttt{ate} method averages predicted contrasts over the supplied
$X$ rows. It cannot recover missing heterogeneity merely by averaging the
same fitted constant repeatedly. A proposed adapter must therefore record
the contrast, treatment scale and reference category, effect-function class,
prediction population and aggregation weights. A marginal effect per unit
of a continuous dose also needs conversion before being interpreted as the
effect of a finite banking action. These fields are part of the scientific
request, not incidental serialization metadata.~\citep{exec19_econml_source}

\subsection{Cross-fitting, numerical conventions and uncertainty have different roles}
\label{exec19:econml-inference}

The inspected final wrapper offers a feature--treatment-residual linear path
and an alternative weighted path. The latter divides the outcome residual
by a signed treatment residual whose magnitude is clipped below at $10^{-5}$,
and uses the original squared treatment residual as a weight. That convention
needs to be preserved in a numerical comparison; it does not supply missing
treatment support. The wrapper should report the convention and diagnostics
without presenting a finite result as evidence that every requested contrast
has adequate overlap.~\citep{exec19_econml_source}

The common learner uses group-aware folds when groups are supplied and a
compatible splitter is used; setting \texttt{cv=1} explicitly disables
cross-fitting. The selected static linear covariance implementation offers
HC0, HC1 and nonrobust choices. It uses a pseudoinverse in the variance
calculation and warns that inference is invalid when the design is rank
deficient. Keeping a borrower's rows in one nuisance fold is useful for
preventing a particular form of training leakage. It does not itself change
the covariance formula into a cluster or calendar estimator. Rank and
sampling warnings must therefore survive the proposed CLI or MCP boundary,
including when the backend returns numerical arrays.
\citep{exec19_econml_source}

A small construction makes the sampling issue explicit. Suppose four customer
score contributions all equal a shared fair shock $Z\in\{-1,1\}$. Each
has mean zero and variance one, but their average equals $Z$ and hence has
variance one. A working calculation that treats the four contributions as
independent gives $1/4$. Group labels in a fitting call cannot remove the
off-diagonal covariance terms. This example concerns a declared common-shock
law; it is not a measured coverage result for EconML.
\citep{exec19_econml_reference}

\subsection{A dynamic panel requires an explicit clock and complete histories}
\label{exec19:econml-panel}

The DynamicDML guide considers treatments assigned over time with all relevant
dynamic confounders observed. Its displayed structural family combines a
state transition linear in previous treatment and state, a potentially
nonlinear assignment policy, and an outcome equation with baseline effect
heterogeneity. The selected implementation targets the final-period outcome,
fits treatment effects backwards through the periods, and uses the first
period's $X$ for heterogeneity. Later $X$ enters as time-varying controls.
These are meaningful design choices for a specified sequential problem; they
do not establish an unrestricted hidden-state servicing or optimal-policy
value guarantee.~\citep{exec19_econml_source}

The code assigns a group's period from its occurrence count in the input.
Its prefit check requires each group to have the same number of rows, but
equal counts do not verify timestamps. For example, the ordered rows
$(a,0),(b,1),(a,1),(b,0)$ contain two dates for each customer. Occurrence
counting nevertheless assigns customer $b$'s date 1 to period 0 and date 0
to period 1. A future adapter must require unique unit--time keys and
canonical chronological histories before making the call. It should also
verify identical unit alignment across period slices, because period-specific
arrays are subsequently combined positionally.
\citep{exec19_econml_source,exec19_econml_reference}

An unbalanced panel raises an error whose text suggests padding the data.
Padding has scientific consequences: an absent contact, an unobserved
repayment and an observed zero are different events. The proposal therefore
requires a declared observation and censoring treatment before any such
conversion. The same principle applies to weights. Although the selected
\texttt{fit} signature accepts sample weights and sample variances, its code
warns that these are unsupported and passes neither into fitting. A requested
weighted population answer must be refused under this profile, rather than
returned as an unweighted result with its warning discarded.
\citep{exec19_econml_source}

The dynamic analytic covariance does include cross-period within-unit score
products. It constructs a block sensitivity matrix and a block score
covariance, then divides the resulting sandwich by the number of groups.
This deserves recognition as more than a rowwise static variance calculation.
It still does not establish covariance qualification for shared calendar
shocks across customers. The proposed request must declare both the history
model and the between-unit sampling law.~\citep{exec19_econml_source}

\subsection{The inspected bootstrap composition can change a history}
\label{exec19:econml-bootstrap}

The source exposes an additional issue that should be resolved before dynamic
bootstrap admission. The common CATE inference mapping selects
\texttt{BootstrapInference}; that object calls \texttt{BootstrapEstimator.fit}.
The latter asks the estimator for strata. DynamicDML returns its group labels.
The bootstrap then draws rows with replacement inside each distinct stratum
and concatenates the draws. Consequently, this composition holds the original
units fixed and resamples their time rows internally. It is not a draw of
whole ordered histories from a population of units. This conclusion follows
from the inspected source call path; no installed backend was run here.
\citep{exec19_econml_source}

For an authored two-period history $(0,1)$, two independent uniform draws from
its rows give
\begin{equation}
 (0,0),\ (0,1),\ (1,0),\ (1,1),
 \qquad P(\hbox{original ordered history})=\frac14.
 \label{exec19:econml-resampling}
\end{equation}
Only $(0,1)$ preserves that history; for two independent unit draws the
probability that both histories survive is $1/16$. Since the selected dynamic
code infers time from occurrence order, duplicated or reversed rows can be
treated as successive periods. This enumeration establishes a contract risk,
not a measured coverage shortfall or a general indictment of the package's
other inference modes.~\citep{exec19_econml_reference}

The selected upstream dynamic tests exercise shapes and inference choices,
including a two-draw bootstrap configuration, and compare selected panel
layouts under automatic analytic inference. Those layouts preserve the
intended chronological information. The selected common bootstrap tests
exercise interval shapes, options and treatment-stratified examples. These
tests were inspected rather than executed. They do not supply the proposed
whole-history conformance evidence. The next prototype should instrument the
pinned resampler with unit, row and time identifiers, preserve the observed
draws, and compare the actual path against an independently specified panel
sampling experiment. A whole-unit bootstrap would itself need an appropriate
independent-unit law; common calendar dependence cannot be repaired by that
substitution without a separate argument.~\citep{exec19_econml_source}

\subsection{Selective reference now, a testable adapter proposal next}
\label{exec19:econml-route}

EconML is a useful reference for flexible effect-function estimation and for
comparison with separately qualified DoubleML profiles. The proposed route
is a narrow Python boundary that preserves targets, histories, warnings and
inference choices. Direct adoption of every estimator and a wholesale port
would enlarge the verification obligation without resolving the scientific
questions above. A transparent finite residual-moment implementation provides
a cheap numerical reference; an estimator with a different longitudinal
contract may be preferable when the selected dynamic structural family is
inappropriate. Refusal remains necessary when identification is unsupported.

The package's root notice contains MIT terms and explicit BSD 3-Clause notices
for specified scikit-learn-derived tree and generalized-random-forest code.
The pinned metadata declares Python at least 3.9, eleven direct dependency
bounds and optional integrations. Its native build requirements include Cython
at least 3.0.11 and below 3.3. These declarations inform feasibility and
maintenance costs; they are not a resolved environment, a complete transitive
notice inventory or proof that a future installation reproduces the selected
source. A prototype needs exact artifacts and applicable notices, with optional
services enabled only when that profile requires them.
\citep{exec19_econml_source}

The first public or synthetic prototype should keep one identified static
target, an explicit effect-function class and a fixed prediction population.
It should reproduce the exact nuisance oracle, test keyed reordering and
treatment-scale conversions, and retain the misspecified constant comparison.
A subsequent dynamic prototype should add complete-history and weight
refusals, trace resampling, and compare declared independent-unit and
common-shock designs. Coverage, interval availability, width and total
failed-run cost should be specified before final evaluation. Such work can
establish whether a narrow implementation merits an expensive bank trial.
The assumed USD amounts above motivate the decision; observed bank prices,
assignment truth and institutional acceptance remain later evidence.
\citep{exec19_econml_reference}


% SOURCE docs/release/2026-10-29/source/docs/writing-repair/2026-10-29/src/aacf2b7f5ffe6d4d.tex CHAR 0
\section{Python effect estimators: population, clipping and the last covariance step}
\label{exec20:pyfx-scope}
A bank evaluating retention offers must first decide whether it wants the
effect for all eligible customers, for customers historically offered the
intervention, or for a risk-defined population facing a future decision.
The same Python method name can hide different averaging, nuisance-fitting
and uncertainty choices. A reusable service must retain these choices in
its request and result. This section compares the already inspected EconML,
DoubleML, CausalML, Delicatessen and statsmodels paths, then develops two
conversion hazards: propensity clipping and the final covariance calculation.
These selected paths do not exhaust Python's effect-estimation ecosystem or
qualify continuous-dose, longitudinal or survival families by association.

\subsection{Five routes with different scientific responsibilities}
\label{exec20:pyfx-routes}
The selected EconML DML path residualizes outcome and treatment and fits an
effect model. Averaging predictions through an ATE interface retains that
model's restrictions; a misspecified constant projection does not become
the population effect merely because its output is called an average.
Its grouping for nuisance splits must also be distinguished from the
sampling law used for covariance. The earlier dynamic discussion shows why
unit histories, encounter order and bootstrap behavior need separate
inspection.~\citep{exec19_econml_source}

The inspected Python DoubleML IRM route makes the score, external nuisance
predictions, folds and repetitions central objects. Its ATE and ATTE choices
are distinct; the R implementation and other model families retain their
own contracts. Earlier matching of seven configurations with fixed identical
nuisance predictions and folds is useful bounded runtime evidence, but does
not qualify new learners, unrestricted dependence or a complete package
distribution. In particular, nuisance vectors must be joined by stable row
keys when another repository prepares them.~\citep{runtime_doubleml_scope}

The selected CausalML T-learner instead fits arm-specific outcomes and reports
contrasts in a defined treatment-label order. An average over prediction rows
depends on the chosen prediction population. Its row bootstrap is not a
borrower-history or common-date bootstrap, and prediction percentiles and
the default ATE interval formula are different objects. An initial binary
treatment profile avoids silently extending that limited review to every
multiarm variance calculation.~\citep{exec14_causalml_primary}

Delicatessen is attractive when the proposal needs transparent stacked
estimating equations and joint nuisance/target inference. The inspected
implementation differentiates the supplied equations and builds a sandwich
from their contribution columns. Those columns encode the sampling units;
the routine does not reconstruct household or calendar dependence from
labels supplied elsewhere. Its allowed pseudoinverse can make a calculation
available without establishing that the target is locally identified.
Analytic small-model derivatives, rank checks and a defensible sampling law
are therefore part of the proposed use.~\citep{v13_sw_delicatessen,v131_sw_delicatessen}

The selected statsmodels \texttt{TreatmentEffect} path supplies five
point-treatment procedures: inverse probability weighting (IPW), regression
adjustment (RA), augmented IPW (AIPW), AIPW with weighted outcome fits, and
IPW regression adjustment. Its documented outcome model is linear OLS.
It combines target, outcome and selection moments for inference where
applicable. This is a useful parametric comparison route, with a different
nuisance-fitting architecture from a cross-fitted flexible learner.
The inspected revision is \texttt{8278e2d218cc} of 6 October 2026; the
following claims concern its selected source, not an unverified later
release.~\citep{v13_sw_statsmodels,exec20_python_effect_source}

For a proposal, the choice is consequently conditional. Use a declared DML
or IRM profile when its score and effect model match the question; retain
arm-specific fits for a separately qualified CATE comparison; use explicit
stacked equations or the parametric statsmodels paths as interpretable
baselines. Reimplementing all five would duplicate model, variance and
failure semantics before a measured need exists. Original engines behind
typed boundaries are the preferred candidates, with executable admission
limited to profiles that later pass their source, environment and behavior
checks. DoWhy, zEpid, causallib and CausalPy remain complementary leads for
target-specific inspection; naming them here grants no capability finding.

\subsection{All customers, treated customers and untreated customers}
\label{exec20:pyfx-population}
Let $A\in\{0,1\}$ denote the historical offer, $X$ baseline information,
$Y(a)$ a contribution outcome under intervention $a$, and
$\tau(X)=\mathbb E[Y(1)-Y(0)\mid X]$. Under consistency, conditional
exchangeability and positivity in the chosen eligible population,
$\tau(X)$ is identified by the corresponding conditional outcome contrast.
Three averages still differ:
\begin{equation}
 \theta_{\rm all}=\mathbb E[\tau(X)],\quad
 \theta_{\rm treated}=\frac{\mathbb E[p(X)\tau(X)]}{\mathbb E[p(X)]},\quad
 \theta_{\rm untreated}=\frac{\mathbb E[(1-p(X))\tau(X)]}{\mathbb E[1-p(X)]},
 \label{exec20:pyfx-targets}
\end{equation}
where $p(X)=\Pr(A=1\mid X)$. The latter two expressions follow by conditioning
on $X$ and using the historical assignment probabilities. They refer to
historical populations; an effect for a future targeting rule requires that
rule's population and information set.

For an exact authored example, suppose low- and high-risk customers are
equally common. Their offer probabilities are $1/4$ and $3/4$, and their
conditional effects are USD 20 and USD 80. Then
$\theta_{\rm all}=50$, $\theta_{\rm treated}=65$ and
$\theta_{\rm untreated}=35$. At a stipulated incremental cost of USD 55,
the corresponding net averages are USD $-5$, USD 10 and USD $-20$.
Selecting the treated-population answer for an all-customer proposal could
therefore reverse the investment story. These are imposed-law calculations,
not observed benefits or a reason to use historical treatment as a future
policy.~\citep{exec20_python_effect_reference}

The inspected statsmodels methods expose these target distinctions through
\texttt{effect\_group}. RA averages fitted potential outcomes over the
selected rows. IPW-RA changes both arm-fitting weights and prediction rows:
for the treated target, for example, treated-arm weights are one and
control-arm weights are propensity odds. AIPW and AIPW-WLS share
target-specific potential-outcome terms. Their joint moments reorder
selection scores to match the grouped outcome rows. An adapter must preserve
these masks and row identities through every nuisance component, rather
than rely on position after independent sorting. Arbitrary user-weight
arrays are outside the initial proposed profile until their point and
inference paths are separately traced.~\citep{exec20_python_effect_source}

\subsection{Clipping an estimated probability changes the score}
\label{exec20:pyfx-clipping}
Consider the ordinary all-population AIPW score with proposed outcome
functions $m_a(X)$ and a working propensity $q(X)$:
\begin{equation}
 \phi_q=m_1(X)-m_0(X)+\frac{A\{Y-m_1(X)\}}{q(X)}
       -\frac{(1-A)\{Y-m_0(X)\}}{1-q(X)}.
 \label{exec20:pyfx-score}
\end{equation}
Write $\mu_a(X)=\mathbb E[Y\mid A=a,X]$, and suppose $0<q(X)<1$.
Taking the expectation conditional on $X$ gives the exact identity
\begin{equation}
 \mathbb E[\phi_q\mid X]-\{\mu_1(X)-\mu_0(X)\}
 =(p-q)\left\{\frac{\mu_1-m_1}{q}
                 +\frac{\mu_0-m_0}{1-q}\right\}.
 \label{exec20:pyfx-bias}
\end{equation}
The right side vanishes if $q=p$, or if both outcome functions equal their
conditional means. Replacing a correct $p$ by a clipped $q$ generally removes
the first route to a zero bias. This is a score identity; estimation error,
sample splitting and asymptotic remainder control need their own argument.

At the inspected source revision, \texttt{ps\_bounds} defaults to
$(0.001,0.999)$ and probabilities are clipped in both point calculations
and selected joint-GMM moment paths. AIPW-WLS also uses probability-derived
outcome weights; the final target calculation still uses the same
population-specific augmented terms. Clipping can reduce numerical
extremes, but it neither establishes positivity in the source population
nor manufactures the rare untreated observations that inference needs.
\citep{exec20_python_effect_source}

To isolate the consequence, take a constant true propensity
$p=1999/2000$, a clipped value $q=999/1000$, conditional means
$\mu_0=20,\mu_1=80$, and deliberately incorrect working outcome functions
$m_0=m_1=0$. The unclipped expected score is 60. The clipped expectation
is $69970/999$, approximately 70.04, with bias $10030/999$.
At a stipulated USD 65 cost, the true mean margin is USD $-5$, whereas
the clipped-score margin is $5035/999$, approximately USD 5.04. Correct
outcome functions would restore the expected score to 60 even with this
clipping. No native package was run, and no claim is made that a fitted
OLS model in this example would select the imposed zero functions.
The example shows what a blanket statement about robustness after clipping
would have to rule out.~\citep{exec20_python_effect_reference}

A proposed service should therefore return the original probability range,
clipping bounds and counts, target weights and overlap diagnostics. If it
trims observations or changes to an overlap-weighted target, that is another
population contract. A sensitivity display across clipping choices is useful
only when the underlying score and target are named; a visually stable
coefficient is not a proof of identification or calibrated uncertainty.

\subsection{Follow inference through to the final moment covariance}
\label{exec20:pyfx-covariance}
The statsmodels class explicitly states that its preliminary outcome
covariance option does not control the final GMM result. The selected five
wrappers call joint GMM without an alternative \texttt{weights\_method};
the default is \texttt{cov}. The final covariance code reads stored fit
options and recomputes moment covariance. HAC is a separate path requiring
a lag choice and kernel; the selected wrapper signatures do not expose that
choice merely by accepting a preliminary robust covariance label.
Thus an HC1 label on an intermediate weighted regression cannot establish
household clustering or common-date uncertainty for the final causal effect.
\citep{exec20_python_effect_source,v131_sw_statsmodels}

Joint nuisance moments can account for estimation of the specified
parametric nuisance model under their sampling and regularity conditions.
They do not justify silently substituting arbitrary machine-learned
predictions, a dependent panel law or a different target score. Record the
full moment vector, derivative matrix, numerical residuals and rank. An
optimizer returning a result establishes a numerical event, not the required
orthogonality, model correctness or sampling approximation.

One small source detail illustrates why that numerical record matters.
The selected covariance routine centers a moment matrix by
\texttt{moms.mean()}, a scalar mean over all entries. Let each row
$m_i\in\mathbb R^k$ be one observation's moments, let
$\bar m=n^{-1}\sum_i m_i$, and set $c=k^{-1}\mathbf1^{\mathsf T}\bar m$.
Denote the second-moment matrices after column and scalar centering by
$C_{\rm col}$ and $C_{\rm scalar}$. Expansion gives
\begin{equation}
 C_{\rm scalar}-C_{\rm col}
   =(\bar m-c\mathbf1)(\bar m-c\mathbf1)^{\mathsf T}.
 \label{exec20:pyfx-centering}
\end{equation}
Cross terms vanish because the column-centered rows sum to zero.
If all moment means are exactly zero, the two constructions agree.
At nonzero residual means they can differ, especially when moment scaling
and an ill-conditioned derivative amplify a numerical discrepancy.
This identity is a qualification question, not a measured software defect
or an allegation of invalid intervals in every fit.
\citep{exec20_python_effect_source}

For the authored moment rows $(1,0)$ and $(3,0)$,
$\bar m=(2,0)$ and $c=1$. The two matrices are
\begin{equation}
 C_{\rm col}=\begin{pmatrix}1&0\\0&0\end{pmatrix},\qquad
 C_{\rm scalar}=\begin{pmatrix}2&-1\\-1&1\end{pmatrix}.
 \label{exec20:pyfx-centering-example}
\end{equation}
The difference is the outer product of $(1,-1)$ with itself. These rows
deliberately do not solve zero-mean moment equations. A conformance test
should therefore include both an exact root and controlled nonzero
residuals, recording the derivative and covariance response rather than
concluding from this example that a particular native estimate is wrong.
\citep{exec20_python_effect_reference}

\subsection{A public prototype that tests meaning before scale}
\label{exec20:pyfx-design}
The first comparison should freeze one binary intervention, units, an
eligible population and independent sampling units. Use the exact target
and clipping laws above, followed by generated heterogeneous outcomes with
known assignment, intentional nuisance misspecification and weak overlap.
Fix row keys, folds, repetitions, predictions and treatment-level order.
For each engine, compare the target expression and retained population
before comparing numbers. Evaluate uncertainty only under its declared law;
then add separately justified clustered or temporal regimes.

The same contract should travel through library, CLI and MCP requests:
target population, treatment contrast, nuisance identities, weighting and
uncertainty settings enter together. The response preserves point estimate,
uncertainty, diagnostics and unavailable states separately. Under an
incorrect model or unsupported dependence law, a successful tool call
must not become a qualified causal claim. Package upgrades reopen the
relevant conversion checks, including row sorting, clipping defaults,
target labels, covariance construction and error handling.

The proposed comparison can measure whether this semantic discipline helps
analysts detect consequential mistakes before an expensive bank experiment.
Setup and dependency maintenance, nuisance training, memory, conversions,
diagnostics and repeated inference are cost categories to measure, not
current prices or savings. Complete distribution provenance and rights,
target-specific native fixtures and broader original-source coverage remain
open. Five reviewed implementations supply concrete choices and tests;
they do not establish a completed ecosystem survey or universal inference.


% SOURCE docs/release/2026-10-29/source/docs/writing-repair/2026-10-29/src/c12ec242576900f2.tex CHAR 0
\section{Java effect engines: uplift, graph-compatible slopes and decision value}
\label{exec20:javafx-scope}

A servicing team may wish to rank customers by the incremental benefit of an
offer, while a treasury analyst may wish to rank possible causes of a funding
response. Both tasks can return a number called an effect, but the corresponding
software contracts differ. H2O-3 supplies a Java implementation of uplift
forests with separate treated and control predictions. Tetrad supplies
graph-compatible regression effects and a stability-ranking construction.
The proposed common library should preserve these useful distinctions rather
than turn every numeric output into a causal confidence interval. The selected
H2O source declares version 3.47.0 and a bleeding-edge profile at the pinned
20 September 2026 revision. Tetrad is inspected at its pinned 22 September 2026
revision. These identities specify the primary text being compared; they do
not certify a released binary or a completed integration.
\citep{exec20_javaeffects_h2o,exec20_javaeffects_tetrad}

\subsection{The target precedes the ranking algorithm}
\label{exec20:javafx-target}

Let $A\in\{0,1\}$ denote a servicing action, $Y(a)$ a binary desirable
response by a declared horizon, and $X$ the information available before the
action. The conditional effect of interest is
\begin{equation}
 \tau(x)=E[Y(1)-Y(0)\mid X=x].
 \label{exec20:javafx-cate}
\end{equation}
For the proposed public comparison, assume consistency, no interference,
conditional exchangeability of the two potential responses and $A$ given $X$,
and positive probability of both actions in the relevant covariate region.
Then $E[Y(a)\mid X=x]=E[Y\mid A=a,X=x]$, so the difference of the
two conditional response means identifies (\ref{exec20:javafx-cate}). This
follows by replacing the conditioning on $A$ under exchangeability and then
using consistency; a treatment-column name alone supplies neither step.
The argument is conditional on these imposed assumptions, not an assertion
that a bank's historical assignments satisfy them.
\citep{exec20_javaeffects_reference}

An exact example makes the cost of omitting assignment information visible.
Suppose half the customers are in each of two pre-action risk groups. Their
action probabilities are $1/5$ and $4/5$, and their response probabilities are
also $1/5$ and $4/5$. Within each group impose identical potential responses,
$Y(1)=Y(0)$, so the causal effect is zero. If the group is omitted,
\begin{equation}
 E[Y\mid A=1]=\frac{17}{25},\qquad
 E[Y\mid A=0]=\frac{8}{25},\qquad
 E[Y\mid A=1]-E[Y\mid A=0]=\frac9{25}.
 \label{exec20:javafx-confounding}
\end{equation}
For example, the first conditional mean is
$[\tfrac12(\tfrac15)^2+\tfrac12(\tfrac45)^2]/(\tfrac12)=17/25$.
With an assumed USD 100 value per incremental response and USD 20 action cost,
the unadjusted contrast suggests USD 16, while the stipulated causal margin is
USD $-20$. This construction does not show that an uplift forest fails when
the confounder is supplied and the model is adequate. It shows why assignment,
covariates, eligibility and target population must accompany a proposed engine.
\citep{exec20_javaeffects_reference}

\subsection{H2O's Java training and scoring contract}
\label{exec20:javafx-uplift}

The selected H2O builder accepts binomial response modeling and requires a
categorical treatment column with two levels, no missing treatment values and
internal extrema zero and one. These checks do not by themselves bind a
business label such as ``contacted'' to the intended treatment code. The
wrapper must save the category domain and the desirable response event.
The selected training path explicitly rejects weights, offsets, a fold column
and a positive built-in cross-validation fold count. A proposed weighted
observational analysis therefore cannot silently pass its weights through this
interface. External held-out evaluation is a separate possible design, subject
to its own sampling and training separation.
\citep{exec20_javaeffects_h2o}

The uplift tree uses a common split structure for both arms and retains two
leaf predictions. For an ordinary leaf with $n_1,n_0$ treated and control
observations and $s_1,s_0$ positive responses, the documented smoothing is
\begin{equation}
 \widehat p_1=\frac{s_1+1}{n_1+2},\qquad
 \widehat p_0=\frac{s_0+1}{n_0+2},\qquad
 \widehat u=\widehat p_1-\widehat p_0.
 \label{exec20:javafx-smoothing}
\end{equation}
The inspected ordinary split code constructs these smoothed rates; the forest
and the selected MOJO scorer average the two probability sums over trees and
return their difference. The three prediction entries are the uplift score,
the treated positive-response prediction and its control counterpart. A
probability difference is not the probability that this particular customer
benefits: that individual statement additionally depends on the joint law of
the two potential responses.
\citep{exec20_javaeffects_h2o}

For ten observations in each arm with eight versus two positive responses,
the raw contrast is $3/5$, but (\ref{exec20:javafx-smoothing}) gives
$3/4-1/4=1/2$. At an assumed USD 100 response value and USD 40 cost, these
two point calculations give USD 20 and USD 10. The difference is an explicit
consequence of smoothing, not empirical evidence of which estimate is better.
If the control arm has no observations, the same formula supplies $1/2$ for
its probability. That numerical value does not identify an unobserved control
response. The result contract should report the missing support alongside any
regularized prediction.
\citep{exec20_javaeffects_reference}

Missing-feature routing and a tree that never splits require additional
qualification. In the selected nonsplitting-root branch, the two leaf values
are taken from the two entries of \texttt{priorClassDist}. The route and
meaning of that inherited object must be traced before treating this branch
as an arm-specific estimate. The selected upstream depth-zero test checks that
a model exists; it does not resolve that semantic question. Other inspected
tests do assert prediction-column names and Java/MOJO scoring comparisons,
which are useful acceptance designs but have not been executed in this work.
\citep{exec20_javaeffects_h2o}

\subsection{A split criterion needs an executable mathematical identity}
\label{exec20:javafx-split}

H2O documents KL, squared Euclidean and chi-squared divergences for uplift
splits. The selected default resolves \texttt{AUTO} to KL, and the histogram
constructor selects the corresponding divergence object. The inspected KL
method contains the expression
\begin{quote}
\ttfamily\small
p * log2(p / q == 0 ? epsilon : q).
\end{quote}
Here $p$ and $q$ abbreviate its two probability arguments. The division and
comparison determine the conditional branch passed to the logarithm; the
expression is not the documented $p\log(p/q)$. For $p=q=1/2$, each selected
literal contribution is $-1/2$, so the two-class node sum is $-1$. The
documented KL sum is zero. The discrepancy persists regardless of the
logarithm base used for the documented zero.
\citep{exec20_javaeffects_h2o,exec20_javaeffects_reference}

This is a bounded source-and-arithmetic finding. It is not a report of an
executed Java failure, a statement about all H2O versions, or a measurement of
the resulting fitted forest. The next public prototype should evaluate the
selected method, normalization and actual split choices on equal probabilities,
zero-support boundaries and asymmetric arm counts. A literal port would carry
the discrepancy; replacing the formula without preserving the original would
create a new algorithm that needs separate review. An alternative divergence
is a candidate to compare, not an already-qualified repair.
\citep{exec20_javaeffects_h2o,exec20_javaeffects_reference}

\subsection{Average predictions and AUUC answer different questions}
\label{exec20:javafx-metrics}

The selected metric builder accumulates predicted uplift values and reports
their averages over all retained rows, treated rows and control rows as ATE,
ATT and ATC. It skips rows with a missing outcome. Thus both the prediction
function and the retained-row population matter; the labels do not by
themselves establish an independently identified or unbiased causal estimate.
For fixed scores $(0.1,0.2,0.3,0.4)$ with the first two rows treated, the
three averages are $1/4,3/20,7/20$. Excluding the fourth row because its
outcome is missing changes the all-row and control averages to $1/5$ and
$3/10$, while the treated average remains $3/20$. A production result must
explain such changes in its population rather than merely label them ATE.
\citep{exec20_javaeffects_h2o,exec20_javaeffects_reference}

For a prefix of customers ordered by a score, let $N_1,N_0$ be arm counts
and $S_1,S_0$ positive-response counts. The selected AUUC implementation uses
\begin{equation}
 L=\frac{S_1}{N_1}-\frac{S_0}{N_0},\qquad
 Q=S_1-S_0\frac{N_1}{N_0},\qquad
 G=(N_1+N_0)L,
 \label{exec20:javafx-curves}
\end{equation}
for lift, qini and gain. These expressions require support in the relevant
arms. Quantile thresholds determine bins; its area convention is a sum of
the cumulative curve value times the bin frequency, divided by $n+1$.
This convention must be fixed in a cross-engine comparison. The selected
implementation replaces a sole NaN curve value with zero and otherwise calls
an interpolation helper for NaNs. Such display processing does not create
the missing assignment information.
\citep{exec20_javaeffects_h2o}

Consider four ordered bins, each containing one treated and one control
customer, with observed response pairs $(1,0),(0,0),(1,1),(0,1)$. Their
cumulative qini values are $1,1,1,0$; the stated rectangle convention gives
$2(1+1+1+0)/9=2/3$, and the endpoint random-reference curve is zero. Scores
$(0.4,0.3,0.2,0.1)$ and their doubled values produce exactly the same order
and these same bins. Their average predictions are $1/4$ and $1/2$. Thus a
matched ranking curve does not establish probability-difference calibration.
This fixed-label arithmetic does not assert any potential-outcome law for
the eight customers. A causal policy comparison needs an independent design,
assignment support, a suitable value estimator and uncertainty for its sampling
unit. Reusing a validation sample repeatedly to choose a threshold also needs
to be included in that evaluation design.
\citep{exec20_javaeffects_reference}

\subsection{Tetrad's effects retain direction, graph ambiguity and failure}
\label{exec20:javafx-graphs}

Tetrad's selected IDA interface works with continuous data and legal partially
directed acyclic graphs. It enumerates admissible local parent choices and
computes regression coefficients using covariance submatrices, with indices
matched by variable names. This preserves a set of candidate linear slopes,
not a general nonlinear heterogeneous response surface. The unchanged IDA
source also underlies the earlier identification discussion; its reuse here
does not add an independent algorithmic coverage result. A selected singular
inverse catch returns zero, while other candidate exceptions can be logged
and omitted. These branches motivate a typed distinction among a graph-implied
zero, a numerical failure and an incomplete candidate collection.
\citep{exec20_javaeffects_tetrad}

For example, if two regressors are identical, a covariance system can determine
only $\beta_X+\beta_Z=1$. The solutions $(0,1),(1,0),(-1,2)$ give the
same fitted combination. Returning zero for the first coefficient does not
show that its causal effect is identified as zero. This exact algebra is a
proposed numerical qualification case; it has not been run through Tetrad.
\citep{exec20_javaeffects_reference}

CStaR extends the comparison from a graph effect set to a magnitude ranking
across subsamples. Its selected code takes the minimum of the absolute
candidate effects. For candidates $\{-2,-1/2\}$, this is $1/2$, whereas
taking the absolute value of the signed minimum gives $2$. The first is the
inspected implementation's intended magnitude summary, but neither operation
retains the sign. At a hypothetical USD 100 value per response unit and USD 20
action cost, the two signed candidates imply margins $-220$ and $-70$.
Misreading the minimum magnitude as a positive benefit gives USD 30. An
integration should retain the signed candidates and their units beside the
ranking summary.
\citep{exec20_javaeffects_tetrad,exec20_javaeffects_reference}

The selected default uses thirty half-samples without replacement; its
bootstrap option also draws half the rows, with replacement. Neither route
alone specifies branch-cluster or shared-calendar resampling. The ranking
uses an inclusive cutoff at the requested bracket times the number of effect
variables. With one effect, bracket one, and three equal positive magnitudes,
all three meet the cutoff. If this happens in every subsample, each has
selection frequency one. That frequency is a property of the construction;
it is not automatically a probability of causation or an interval's coverage.
The displayed PCER calculation uses the displayed rank as its $q$ argument.
Its presence does not verify the theoretical assumptions behind an error bound
for the proposed data. The source also states that required and forbidden
edge knowledge is not respected by this class. A business constraint must
therefore be preserved through a verified route or explicitly refused.
\citep{exec20_javaeffects_tetrad,exec20_javaeffects_reference}

\subsection{A proposed adoption route and a public qualification experiment}
\label{exec20:javafx-adoption}

H2O is a candidate for an original Java training worker with a separate MOJO
scoring service; Tetrad is a candidate for graph-effect and stability-ranking
reference services. The two should expose different result types through the
library, CLI and MCP surfaces. A request should bind row and variable identities,
action and positive-event mappings, horizon, eligible population, graph type,
original units, sampling law, source revision and executable artifact hash.
A response should separate prediction, causal estimand, graph candidates,
ranking metric, numerical status and qualified uncertainty. None of these
proposed Java services has been implemented or admitted by this source review.
\citep{exec20_javaeffects_h2o,exec20_javaeffects_tetrad}

Preserving the original engines is preferable as the first comparison because
it makes a source-to-output discrepancy observable before maintaining a port.
H2O's training and generated-model modules have different dependencies; the
generated-model build's compatibility setting is not a statement about every
training deployment. Its selected MOJO subclasses both declare format 1.40,
but use different threshold-key strings in their writer and reader. Superclass
and default handling must be traced and a full archive roundtrip tested before
concluding whether this affects scoring or metadata. The root Apache 2.0
declaration and Tetrad's GPLv3-or-later source headers are relevant inputs to
reuse review; neither a process boundary nor a source notice resolves the
complete distribution and dependency obligations.
\citep{exec20_javaeffects_h2o,exec20_javaeffects_tetrad}

The first public experiment should therefore use frozen synthetic mechanisms
with known action effects, followed by appropriately documented public data.
The synthetic suite should include the confounded zero-effect law, missing arm
support, constant/root-only trees, missing features, category and column
permutations, the equal-probability split witness, signed graph effects,
singular covariance and ranking ties. Separate checks should match engine and
MOJO predictions and preserve keyed data through each transport. After that
semantic comparison, a fresh held-out study can evaluate calibration, ranking,
policy value and repeated-fit costs under declared independent or dependent
sampling laws. No bank experiment is required to establish whether these
public prototype steps merit the next investment.
\citep{exec20_javaeffects_reference}

This review supplies two concrete Java profiles rather than an absence finding
from an empty search. Saved complementary repository leads and search remainder
still need original-source inspection. Foundational papers named in upstream
comments need genuine local PDFs and technical review before their theorems
are attributed here. Those tasks can proceed alongside the pinned public
prototype design; their openness does not justify converting a score, a
frequency or a successful API call into an unsupported causal guarantee.
\citep{exec20_javaeffects_h2o,exec20_javaeffects_tetrad}


% SOURCE docs/release/2026-10-29/source/docs/writing-repair/2026-10-29/src/d587fa2f22fbb93b.tex CHAR 0
\section{Native causal forests: preserve the target before accelerating it}
\label{exec21:cxx-scope}
A servicing team considering different offers for different risk groups needs more than an accurate forecast of repayment. It needs a contrast between actions, a declared population over which to average that contrast, and uncertainty that corresponds to the decision. A C++ implementation can reduce computation while changing any of these meanings at a language boundary. The useful comparison therefore starts with the statistical object and follows it through wrappers, native learning, stored draws and portfolio aggregation. Four inspected lineages provide complementary designs: generalized random forests, \texttt{bcf}, XBCF and flexBCF. Their common use of trees does not make their targets, priors or intervals interchangeable.~\citep{exec21_cxx_grf,exec21_cxx_bcf,exec21_cxx_xbcf,exec21_cxx_flexbcf}

The comparison binds grf 2.6.1 to its 30 April 2026 commit, bcf 2.0.2 to 26 September 2023, XBCF's declared 0.1.2 source to 25 June 2024, and flexBCF 1.0.3 to 1 December 2025. These are inspected versions preceding the literature cutoff, not assertions about the newest release. The native grf boundary analysis developed elsewhere in this monograph remains applicable: the initial proposal favors a qualified R process around its native core, rather than a new direct embedding or port. This comparison extends the alternatives and semantic contracts; it supplies no runtime qualification.~\citep{exec21_cxx_grf,exec21_cxx_bcf,exec21_cxx_xbcf,exec21_cxx_flexbcf}

\subsection{A conditional effect and a dose slope answer different questions}
\label{exec21:cxx-target}
For a binary action $A$, let $Y(a)$ be the potential outcome under action $a$ and $X$ permitted baseline information. Under consistency, no interference, conditional ignorability and overlap, the conditional average effect is
\begin{equation}
 \tau(x)=E\{Y(1)-Y(0)\mid X=x\}.
 \label{exec21:cxx-cate}
\end{equation}
The inspected grf wrapper describes this binary target, estimates or accepts aligned outcome and treatment nuisance predictions, and residualizes before forest fitting. Honesty separates split selection from leaf estimation. The native trainer checks whether the sample and honesty fractions leave observations in both parts; its tree trainer splits sampled clusters and repopulates estimation leaves. Computational separation does not establish ignorability or validate a treatment clock.~\citep{exec21_cxx_grf}

For continuous treatment the wrapper documents a conditional covariance-to-variance slope,
\begin{equation}
 s(x)=\frac{\operatorname{Cov}(Y,A\mid X=x)}{\operatorname{Var}(A\mid X=x)}.
 \label{exec21:cxx-slope}
\end{equation}
It cannot be substituted automatically for $E\{Y(a_1)-Y(a_0)\mid X=x\}$. In an authored example, randomize $A$ uniformly over $\{-1,0,1\}$ and set $Y(a)=a^2$. The covariance numerator is zero by symmetry, while the action contrast from zero to one is one. Neither confounding nor a computational failure is needed for the discrepancy: the requested objects differ. A bank seeking the response to a specific change in fee or service intensity must request that dose pair or justify a response model under which the slope represents it. This is an exact semantic witness, not a run of grf.~\citep{exec21_cxx_grf,exec21_cxx_reference}

The native prediction strategy solves a weighted local moment ratio. Its instrumental implementation divides an outcome--instrument covariance expression by a treatment--instrument expression. A near-zero denominator is a scientific and numerical warning. The inspected variance code groups trees, skips groups with empty leaves and corrects between- and within-group quantities. Its objective-Bayes correction avoids negative variance estimates; that label does not make its output draws from a Bayesian causal forest. An adapter must retain effective group counts, denominator diagnostics and unsupported variance status. Full upstream guard behavior requires a qualified native trial.~\citep{exec21_cxx_grf}

\subsection{Why regularizing prediction can distort a treatment contrast}
\label{exec21:cxx-prior}
The genuine inspected BCF original is the 33-page arXiv version 4 of 13 November 2019 by Hahn, Murray and Carvalho; it is not asserted to be byte-identical to the later journal article. Under its stated assumptions the conditional mean is modeled as
\begin{equation}
 E(Y\mid X=x,A=a)=\mu\{x,\widehat\pi(x)\}+a\tau(x),
 \label{exec21:cxx-bcf-model}
\end{equation}
with separate tree priors for the prognostic and treatment functions. The fitted propensity $\widehat\pi$ informs the response-surface prior. This addresses regularization-induced confounding: shrinkage of the prognostic function can leave treatment to explain outcome variation associated with selection. A propensity feature does not imply an orthogonal estimating equation or double robustness. The paper explicitly disclaims double robustness for its natural Bayesian estimates and discusses sensitivity to alternative propensity estimates.~\citep{exec21_cxx_hahn}

A project derivation isolates the mechanism. Let $X$ be equally likely to be $-1$ or $1$, with treatment probabilities $1/4$ and $3/4$ respectively. Set $Y(a)=2X$ for both actions, so the true causal effect is zero and overlap holds. Write $Z=A-1/2$. Minimize
\begin{equation}
 E\{(Y-\beta X-tZ)^2\}+\lambda\beta^2,
 \qquad \lambda=3/4.
 \label{exec21:cxx-shrink-objective}
\end{equation}
Only the prognostic coefficient is penalized. Here $E(X^2)=1$, $E(Z^2)=1/4$ and $E(XZ)=1/4$. The treatment normal equation gives $t=2-\beta$; the prognostic equation then gives
\begin{equation}
 \beta=\frac{2(3/4)}{3/4+\lambda}=1,
 \qquad t=1.
 \label{exec21:cxx-shrink-result}
\end{equation}
The apparent treatment effect is produced by the stipulated shrinkage. The exact reference verifies both normal equations. This motivates studying prior design, treatment selection and nuisance learning jointly. It is not a theorem that every regularized learner fails or that BCF necessarily repairs every selection mechanism.~\citep{exec21_cxx_reference,exec21_cxx_hahn}

The bcf wrapper exposes separate prognostic and moderation matrices, supplied propensities and their placement in either or both ensembles. Prior-scale choices and outcome standardization belong to the fitted model. A port that copies trees but drops a scale prior changes the procedure. flexBCF explicitly distinguishes its default model from bcf in scale priors, treatment-recoding behavior and relative regularization of the two ensembles. A claim of speed or flexibility therefore cannot establish identical inference. A prototype comparison must freeze each configuration and compare a common target.~\citep{exec21_cxx_bcf,exec21_cxx_flexbcf}

\subsection{Three meanings of weight must remain separate}
\label{exec21:cxx-weights}
The bcf argument \texttt{w} represents likelihood precision: the conditional variance is $\sigma^2/w_i$. The native source uses these weights in tree updates and residual sums. They do not define the population for a reported effect. For a declared cohort $\mathcal I$ and target weights $q_i$, a draw-wise average is
\begin{equation}
 T^{(s)}=\frac{\sum_{i\in\mathcal I}q_i\tau_i^{(s)}}{\sum_{i\in\mathcal I}q_i}.
 \label{exec21:cxx-aggregate}
\end{equation}
Likelihood precision, target weights and cohort selection require different fields. The target can concern eligible customers, treated customers or a transported population; its meaning cannot be recovered from an unlabeled vector.~\citep{exec21_cxx_bcf,exec21_cxx_flexbcf}

For two equally represented groups with effects $-2$ and $6$, the population effect is $2$. Reusing precisions $1$ and $9$ as target weights gives $26/5=5.2$. If these effects are stipulated dollar gains and the offer costs USD 3 per customer, the population margin is USD -1 while the reweighted margin is USD 2.20. Better observation precision does not authorize replacing the eligible customer mix.~\citep{exec21_cxx_reference}

In grf, equalizing cluster weights changes training samples and average-effect weights. The documented training rule draws equal numbers of units from sampled clusters, using the smallest cluster size. This option is incompatible with supplying the ordinary sample-weight vector in that interface. Balancing branches instead of customers changes both the target and parts of learning; it cannot be implemented by relabeling a final average while leaving an incompatible training run untouched.~\citep{exec21_cxx_grf}

\subsection{Joint draws, scale restoration and persistence}
\label{exec21:cxx-draws}
A Bayesian portfolio interval averages customer effects within a common posterior draw before summarizing portfolio draws. Averaging individual interval endpoints generally differs. Consider two equally weighted customers with joint draws $(-2,6)$ and $(6,-2)$. Each customer's marginal support is $[-2,6]$, but every portfolio draw is exactly $2$. Averaging marginal endpoints discards the perfect negative dependence. At the assumed USD 3 cost every portfolio margin is USD -1. Two artificial draws illustrate algebra; they establish neither convergence nor calibration.~\citep{exec21_cxx_reference}

The inspected flexBCF native SATT routine evaluates each stored ensemble on supplied rows, averages within the draw and then calculates quantiles. Its weighted companion divides by the supplied weight sum. The routine name does not select treated rows: cohort selection must already be reflected in its inputs. Following the exported caller at the same revision reveals a more specific problem. In \texttt{average\_tree\_fits}, the treatment branch selects the treatment trees but passes \texttt{treat = FALSE}. The native routine consequently applies prognostic restoration, multiplying by the outcome standard deviation and adding the outcome mean. By contrast, \texttt{get\_tree\_fits} passes \texttt{TRUE} for treatment trees and its native predictor restores only the scale. The supplied predictors reach the corresponding evaluation fields in either branch; the consequential difference here is the additive outcome mean.~\citep{exec21_cxx_flexbcf,rel16flexcallers}

Let $g_i^{(s)}$ denote a standardized treatment-tree evaluation, let $q_i$ be nonnegative target weights with positive sum, and let $\bar Y$ and $s_Y$ be the training outcome mean and standard deviation. The intended and incorrectly restored portfolio draws are
\begin{equation}
 T^{(s)}=s_Y\frac{\sum_i q_i g_i^{(s)}}{\sum_i q_i},
 \qquad \widetilde T^{(s)}=T^{(s)}+\bar Y.
 \label{rel16:flex-scale}
\end{equation}
This difference shifts the posterior mean and every translation-equivariant posterior quantile by the same amount. Consider an authored draw for twenty equally weighted customers: ten standardized evaluations are $1/5$ and ten are $2/5$. With an outcome mean of USD 100 and standard deviation of USD 10, the intended average gain is USD 3; prognostic restoration produces USD 103. At an assumed USD 5 action cost, the corresponding margins are USD $-2$ and USD 98. The two routes imply opposite decisions despite using the same tree evaluations. These are stipulated exact calculations, not output from an executed package or evidence about an actual bank.~\citep{rel16flexcallers,rel16flexwitnesses}

A useful qualification test follows from the meaning of an effect. Adding a common constant to every observed outcome leaves the centered, standardized training outcomes unchanged and should leave the action contrast unchanged, for the same fitted draws and settings. An unwanted additive outcome mean instead transfers that constant into the reported effect. The adapter should test this invariance and compare draw-wise weighted averages from the individual-evaluation path with the summary path after repairing the branch. Simply subtracting the outcome mean downstream would leave the original caller defect and other boundary contracts unresolved.~\citep{rel16flexcallers,rel16flexwitnesses}

The caller rejects a single-row prediction and negative weights, but the inspected caller and weighted routine do not establish a complete contract for finite aligned weights with positive sum or nonempty retained draws. The native routine accesses the first stored draw and divides by the supplied weight sum. A further small-cohort hazard comes from the default progress interval: for two through nine rows, its floor-of-one-tenth rule gives zero, which reaches an integer remainder operation even when verbose output is disabled. The individual-evaluation caller similarly derives its default from the number of draws. These source paths require explicit positive progress intervals, valid draw counts and input checks before native testing. They are findings about the pinned implementation, not a claim that every entrypoint fails or that a particular exception has been observed.~\citep{rel16flexcallers,rel16flexwitnesses}

Scale parameters must remain aligned by draw. In inspected bcf prediction an ensemble output $g_i^{(s)}$ contributes
\begin{equation}
 \tau_i^{(s)}=s_Y\{b_1^{(s)}-b_0^{(s)}\}g_i^{(s)},
 \label{exec21:cxx-restore}
\end{equation}
where $s_Y$ is the outcome scale. A treatment contrast receives no added outcome mean. XBCF also restores a treatment draw using outcome scale and its treatment-scale difference. If raw draws are $1,2$, $s_Y=10$, and scale differences are $1/2,2$, restored draws are $5,40$ with mean $22.5$. Omitting the contrast factor gives mean $15$. At an assumed USD 20 action cost these imply USD 2.50 and USD -5 margins. Multiplying mean raw output by mean scale also loses their posterior dependence.~\citep{exec21_cxx_bcf,exec21_cxx_xbcf,exec21_cxx_reference}

bcf documents draws by units, whereas selected XBCF outputs use units by iterations. Transport needs unit, chain and retained-draw identifiers rather than a guessed transpose. bcf's wrapper serializes tree samples for later prediction; suppressing file output prevents that path. Its selected prediction wrapper requires both treatment labels, a restriction to reconcile for a one-cohort request. The chain executor sets one core while retaining within-chain thread controls. These are persistence and execution questions for an adapter, not measured speed or scalability.~\citep{exec21_cxx_bcf,exec21_cxx_xbcf}

\subsection{Static hazards define focused prototype tests}
\label{exec21:cxx-hazards}
In pinned XBCF's combined R predictor, the default propensity branch refers to \texttt{z}, which is not a formal argument. The function assigns predictions to \texttt{obj1} and \texttt{obj2}, then reads \texttt{obj} before its local assignment at the end. The separate treatment-draw predictor uses its own local object. These are particular source paths requiring isolated reconciliation, not failures observed by executing the package. An enclosing environment, later release or different entrypoint must not silently replace the inspected route.~\citep{exec21_cxx_xbcf}

In the selected Python wrapper, the zero-outcome-variance branch constructs \texttt{ValueError} without raising it, and adjusted fitted draws are assigned inside the standardization branch. Prediction has separate restoration logic. Tests should cover constant outcomes, both standardization settings, fit-versus-predict parity and return formats. The native loop alternates prognostic and treatment updates and updates scales; stored output must be checked against wrapper order. Source resemblance does not prove equivalent R and Python interfaces.~\citep{exec21_cxx_xbcf}

The selected flexBCF trainer derives the control count from prognostic and treatment design sizes and offsets treatment updates accordingly. Its public wrapper transposes the supplied matrices and passes the treatment indicator without rearranging rows. The same-revision manual explicitly requires controls before treated subjects; its worked construction supplies prognostic predictors for all subjects but treatment predictors only for treated subjects. The native traversal and residual offsets follow that arrangement. An adapter must apply one common permutation to outcomes, treatment, identifiers and prognostic predictors, then form the treated-only design in the corresponding order. Sorting just the treatment vector would corrupt the correspondence. For subjects $(a,b,c,d)$ with labels $(1,0,1,0)$, a stable valid order is $(b,d,a,c)$ and the treatment design contains $(a,c)$. This identity check is an authored fixture, not a fitted forest.~\citep{exec21_cxx_flexbcf,rel16flexcallers,rel16flexwitnesses}

The manual also requires continuous predictors scaled to $[-1,1]$ and zero-indexed categorical levels. The inspected parsers check selected matrix and level-list dimensions; they do not establish validation of every observed category, outcome scale or treatment-row correspondence. The training wrapper standardizes outcomes before the native call, so a constant outcome needs an explicit supported disposition. These checks should be retained alongside the documented input order and scale conventions. Assuming the same row layout across all four libraries would hide a consequential boundary error.~\citep{exec21_cxx_flexbcf,rel16flexcallers}

\subsection{Historical performance supports an experiment}
\label{exec21:cxx-evidence}
The BCF paper offers evidence for its design and relevant limitations. Its targeted-selection example improves bias and coverage over ordinary BART under that simulation; broader experiments do not give uniformly nominal coverage. In the retrospective ACIC 2016 comparison, Table 4 reports nominal 95\% interval coverage of 82\% for BCF, 88\% for propensity-score BART, 81\% for BART and 58\% for the compared causal random forest. BCF's mean interval length is 0.026 versus 0.055 for that forest. Shorter intervals and lower estimation error coexist with undercoverage. The paper discusses this limitation and supplies no general frequentist coverage theorem for its credible intervals.~\citep{exec21_cxx_hahn}

These are historical reported results, not a project reproduction or ranking of today's implementations. They motivate public experiments reporting effect error, coverage, width, failures and decision consequences together. Known counterfactual generators can vary overlap, targeted selection, effect sparsity, response nonlinearity, propensity misspecification and dependence. Public benchmarks provide complementary designs whose versions, targets and access conditions must be fixed. Failed fits and invalid requests stay in the denominator. Posterior uncertainty, Monte Carlo error and repeated-sampling coverage must remain separate: more draws do not repair a wrong causal target or unmeasured confounding.~\citep{exec21_cxx_hahn,exec21_cxx_reference}

\subsection{A scalable longitudinal forest still needs an identification model}
\label{rel17:flex-paper}
The original flexBCF study by \citet{rel17flexpaper} supplies a consequential comparison between computational feasibility and scientific adequacy. The inspected original is arXiv version 2, posted on 11 May 2023 and carrying a 15 May 2023 title-page date. It is distinct from both a fresh replication and an audit of the selected 2025 package. The authors analyze the 2022 ACIC challenge: 3,400 synthetic datasets, comprising 200 replications of each of seventeen mechanisms, with treatment assigned to medical practices after two pre-treatment years. For banking, the analogous question is the effect of a branch-level intervention on customers served by treated branches. It requires a branch assignment clock, a specified customer population and a defense of the comparison across branches.~\citep{rel17flexpaper}

Their identification argument uses consistency and no anticipation, overlap, and conditional independence of untreated outcome changes and assignment given baseline covariates. Define $Z$ as branch assignment, $X$ as permitted baseline information, $s$ as a pre-intervention date and $t$ as a post-intervention date. Let $m_{zt}(x)=E(Y_t\mid Z=z,X=x)$. The corresponding conditional effect among treated units is
\begin{equation}
 \tau_{\rm T}(x,t)
 =m_{1t}(x)-m_{0t}(x)-\{m_{1s}(x)-m_{0s}(x)\}.
 \label{rel17:flex-did}
\end{equation}
The paper proves this difference-in-differences identity in its Appendix S1. Its stated conditional independence is stronger than the equality of conditional untreated mean changes needed for this mean identity. A portfolio ATT must further average over the declared treated population; an unqualified all-customer average is a different target. Neither a fitted propensity feature nor a large number of customer rows verifies the identifying assumptions.~\citep{rel17flexpaper}

The estimation model adds a separate restriction. It represents the outcome mean by a common prognostic function of covariates and time, plus a treatment function activated only after intervention. Before intervention, that specification forces treated and untreated groups with the same included covariates to have the same expected outcome level. Parallel untreated trends alone do not impose equal baseline levels. To see the distinction, consider equally likely branch groups with no varying observed covariates. Their pre-intervention monetary means are USD 10 and USD 14. Untreated outcomes rise by USD 3 in both groups; the second group receives an intervention adding USD 2. Observed post-intervention means are therefore USD 13 and USD 19. The post-period gap is USD 6, but subtracting the USD 4 pre-period gap identifies the USD 2 effect. At an assumed USD 3 per-customer intervention cost, the identified margin is USD $-1$, whereas treating the post-period gap as the effect yields USD 3. This exact law satisfies overlap and parallel untreated increments while violating equal baseline levels. It does not assert that a particular fitted forest necessarily returns the post-period gap.~\citep{rel17flexpaper,rel17flexwitness}

The challenge analysis also conditions the treatment component on some post-treatment practice-composition summaries. The authors explicitly assume that intervention does not change attrition and hence practice composition. A bank must examine that assumption when assistance changes account closure, migration or the customers remaining at a branch. If composition responds to intervention, conditioning on its later value can change the causal question or introduce selection. A serious alternative is to model changes under a defended trend restriction, with a fixed population and dependence-aware evaluation; it is not enough to feed later composition into a flexible learner and retain the original target label.~\citep{rel17flexpaper}

Computational improvement did not remove the inferential limitation in the reported challenge. For nominal 90\% overall-ATT credible intervals, the original beneficiary-level submission had mean coverage 52.15\% across the seventeen mechanisms; the practice-level analysis had mean coverage 83.68\%. Its mean RMSE rose from 20.54 to 23.04 when aggregated to practices. The paper reports that flexible propensity estimation often increased interval width and coverage, while point-estimation differences were modest. These historical summaries are attributed results, not independent reconstructions, and they do not establish a universal preference for aggregation, a propensity learner or a sparsity prior. The authors discuss outcome-model misspecification, omitted pre-treatment outcomes or trends, and limited sensitivity analysis as explanations and limits.~\citep{rel17flexpaper}

The study motivates examining tree persistence, categorical partitioning and memory allocation, but its runtime narrative supplies no matched bank total-cost estimate. Nor does its empirical performance diagnose the later pinned caller defect above: these concern different versioned evidence. For a bank prototype, the useful comparison separates the identifying trend restriction, additional response-model restrictions, customer and branch sampling units, composition assumptions, and the reporting path. Faster computation merits adoption only when the resulting estimand and uncertainty remain appropriate to the decision. A failed calibration or an unsupported composition assumption should reverse an otherwise attractive computational recommendation.~\citep{rel17flexpaper,rel16flexcallers}

\subsection{Integration choices and the remaining census}
\label{exec21:cxx-integration}
For the proposal, grf remains a candidate through its analyzed R/native boundary. bcf is a candidate for a process-based Bayesian adapter with explicit priors and persisted trees. XBCF warrants a performance comparison after selected wrapper paths are reconciled. flexBCF is a distinct alternative needing prior, coding, training-order and cohort contracts; its inspected public treatment-summary path must be repaired and qualified before it is enabled. Porting would transfer responsibility for random numbers, tree updates, priors, draw alignment and persistence. The initial approach should preserve each native model behind a typed adapter and qualify a narrow target before considering a port. These are reasoned proposal choices, not completed reuse admissions.~\citep{exec21_cxx_grf,exec21_cxx_bcf,exec21_cxx_xbcf,exec21_cxx_flexbcf}

Inspected declarations are GPL-family terms for grf, bcf and flexBCF and Apache 2.0 for XBCF. They identify a rights review rather than completing one: dependencies, notices, distribution and embedding remain to be examined. A separate process does not automatically settle licensing obligations. Bounded repository discovery also found SparseBCF, time-series, multivariate, survival and other candidates with CRAN mirrors. Names are leads; lineage, cutoff availability, originals and capabilities remain uninspected. Mirrors do not count as independent implementations.~\citep{exec21_cxx_grf,exec21_cxx_bcf,exec21_cxx_xbcf,exec21_cxx_flexbcf}

Library, CLI and MCP requests should carry action contrast or slope, cohort, clock, source version, nuisance provenance, weight roles, uncertainty type and refusal reasons. An agent can help formulate requests and inspect diagnostics, but cannot relabel posterior quantiles as calibrated confidence limits. These exact witnesses justify inexpensive public prototype tests before costly bank trials. Broader coverage, method originals, rights, native parity and calibrated decision performance remain prerequisites for a stronger adoption claim.


% SOURCE docs/release/2026-10-29/source/docs/writing-repair/2026-10-29/src/3379d7ea6a8fb354.tex CHAR 0
\section{Reusing the R forest package: nuisance provenance and population choice}
\label{exec21:grfr-scope}
A bank considering a retention offer needs an effect for the customers it intends to contact. The R package grf is a useful reference candidate because its wrapper connects nuisance fitting, forest prediction and average-effect estimation. That wrapper is part of the statistical method. A proposed CLI or MCP interface must preserve its population and score construction when exposing the underlying C++ engine. The R package and native core share one implementation lineage; their different deployment routes do not supply independent corroboration. We select the original R package as a bounded reference candidate, with an isolated worker considered for a future prototype. The earlier native-boundary analysis explains the ownership and persistence obligations; here the additional concern is the meaning of supplied nuisance predictions and target options.~\citep{exec21_grf_reuse,exec13_grf_memory}

The selected source declares version 2.6.1 at the 30 April 2026 commit. Its requirements include R 3.5.0 or later, C++17, GNU make, Rcpp and RcppEigen, with imports including DiceKriging, lmtest, Matrix and sandwich. A named maintainer and issue tracker establish a maintenance contact route, not a service commitment. The DESCRIPTION declares GPL-3; file-level terms and dependencies need to accompany any distribution decision. The inspected license distinguishes running a program from conveying it. An isolated process is an engineering boundary and does not by itself settle the licensing of a distributed system. These observations justify retaining the original as a reference without claiming a qualified installed environment.~\citep{exec21_grf_reuse}

\subsection{What the nuisance inputs must mean}
\label{exec21:grfr-nuisance}
Let $A\in\{0,1\}$ denote the offer, $X$ information available before assignment, and $Y$ the outcome over a fixed horizon. Define $e(x)=P(A=1\mid X=x)$ and $\mu_a(x)=E[Y\mid A=a,X=x]$. Consistency, no interference, conditional exchangeability and positivity give $\tau(x)=\mu_1(x)-\mu_0(x)$ the causal interpretation. These are substantive assumptions about assignment and observation; fitting the forest does not verify them.

The wrapper's outcome nuisance is the marginal regression $m(x)=E[Y\mid X=x]=(1-e(x))\mu_0(x)+e(x)\mu_1(x)$. Missing nuisance vectors are estimated by auxiliary regression forests. Supplied scalar values are replicated, and supplied vectors are checked for length. Such checks cannot establish that a prediction uses only permitted information or that it was estimated without access to the evaluation outcome. The interface therefore needs the feature clock, row identity and fitting provenance as well as the numbers. Initially, the proposed worker retains the package's ordinary nuisance path or an explicitly known randomized propensity; externally fitted nuisance vectors require a separate qualification profile.~\citep{exec21_grf_reuse}

Write $m^\star,e^\star,t^\star$ for fixed proposed nuisance functions, with $0<e^\star<1$. The binary score reconstructed by the source is
\begin{equation}
 \Gamma=t^\star+\frac{A-e^\star}{e^\star(1-e^\star)}
       \{Y-m^\star-t^\star(A-e^\star)\}.
 \label{exec21:grfr-score}
\end{equation}
Suppressing $x$, define $u_0^\star=m^\star-e^\star t^\star$ and
$u_1^\star=m^\star+(1-e^\star)t^\star$. Substituting these definitions yields
$\Gamma=u_1^\star-u_0^\star+A(Y-u_1^\star)/e^\star
-(1-A)(Y-u_0^\star)/(1-e^\star)$. For a fresh observation, conditional on these fixed functions, direct expectation gives the project identity
\begin{equation}
 E[\Gamma\mid X]-\tau
 =(e^\star-e)\left\{\frac{u_1^\star-\mu_1}{e^\star}
                   +\frac{u_0^\star-\mu_0}{1-e^\star}\right\}.
 \label{exec21:grfr-bias}
\end{equation}
Indeed, the two residual terms have conditional expectations
$e(\mu_1-u_1^\star)/e^\star$ and
$(1-e)(\mu_0-u_0^\star)/(1-e^\star)$, and collecting errors gives the identity. The conditional bias vanishes when the propensity is correct or when both reconstructed treatment-specific regressions are correct. Correctness of $m^\star$ alone is a different condition. Estimation rates, sample reuse and dependence still govern whether a fitted procedure has valid uncertainty; this fixed-function calculation supplies no finite-sample guarantee.~\citep{exec21_grf_reuse,exec21_grf_reference}

Consider a deliberately simple law with no varying $X$: $P(A=1)=1/4$, $Y(0)=0$ and $Y(1)=2$. The true marginal outcome mean is $1/2$ and the effect is two. Supply the correct values $m^\star=1/2$ and $t^\star=2$, but use $e^\star=1/2$. The reconstructed regressions become $-1/2$ and $3/2$, although the true regressions are zero and two. Equation~\eqref{exec21:grfr-score} gives scores one for controls and three for treated observations, so their population mean is $3/2$. Even a correct marginal outcome regression and correct conditional-effect input do not restore the intended mean in this example.

For comparison, the correct propensity $e^\star=1/4$ with deliberately wrong $m^\star=3$ and $t^\star=0$ produces scores four and minus four, whose mean is two. Alternatively, keeping the wrong propensity $1/2$ but setting $m^\star=1$ and $t^\star=2$ reconstructs the correct outcome surfaces and makes both scores two. Exact rational enumeration checks these cases and the bias identity on 108 combinations. These are supplied-input witnesses, not claims that the fitted grf algorithm generates these nuisance values. Their practical purpose is to prevent an adapter from equating a correct outcome model with an unlabeled marginal prediction vector.~\citep{exec21_grf_reference}

\subsection{An overlap recommendation changes the target}
\label{exec21:grfr-target}
The function \texttt{average\_treatment\_effect} exposes all, treated, control and overlap populations. Under the preceding causal assumptions, their common representation is
\begin{equation}
 \theta_q=\frac{E[q(X)\tau(X)]}{E[q(X)]},\qquad
 q(X)\in\{1,e(X),1-e(X),e(X)(1-e(X))\}.
 \label{exec21:grfr-populations}
\end{equation}
All denominators must be positive. The source warns when estimated propensities become extreme and suggests alternative targets or filtering. A warning does not prove a failure of identification, establish a universal threshold or authorize changing the requested population. A business owner deciding about all eligible customers must explicitly reconsider the question before receiving an effect for customers with greater assignment overlap.~\citep{exec21_grf_reuse}

In a two-group hypothetical bank, let each group contain half the eligible population. Treatment propensities are $1/2$ and $1/10$, and group effects are USD $-4$ and USD 12. The all-customer effect is USD 4. The treated and control population effects are respectively USD $-4/3$ and USD $44/7$. Overlap weights are $1/4$ and $9/100$, giving
\begin{equation}
 \theta_{\rm overlap}
 =\frac{(1/4)(-4)+(9/100)12}{1/4+9/100}
 =\frac{4}{17}.
 \label{exec21:grfr-overlap}
\end{equation}
With an assumed offer cost of USD 2 per recipient, the all-customer margin is USD 2, while the overlap-population margin is USD $-30/17$, approximately USD $-1.76$. Both calculations are correct for their stated populations. Applying the second margin to the original customer population would be a change of decision target, not merely a more stable computation of the same quantity. These dollar figures are reproducible assumptions and do not measure a bank's return.~\citep{exec21_grf_reference}

The implementation also makes the method label consequential. For the all-population AIPW path, it aggregates scores and computes centered cluster sums for the standard error. The overlap path instead regresses the outcome residual on the treatment residual, with an intercept and observation weights, and uses a sandwich variance. Inside that branch the AIPW/TMLE method choice is ignored, although the earlier TMLE guard rejects clustered or sample-weighted inputs. A wrapper should report the actual residual-regression method and preserve the combination restrictions. It should not label an overlap result as an all-population augmented estimate merely because the request said AIPW.~\citep{exec21_grf_reuse}

The population logic of that residual regression can be checked without executing the package. With correct $m$ and $e$, the treatment residual has conditional mean zero, and
$E[(A-e)(Y-m)\mid X]=e(1-e)\tau(X)$ while
$E[(A-e)^2\mid X]=e(1-e)$. The population regression coefficient therefore equals the overlap effect in Equation~\eqref{exec21:grfr-populations}, provided its denominator is nonzero. This argument explains the target. It does not prove that estimated residuals, a particular sandwich correction or a small sample deliver calibrated intervals.

\subsection{Qualification follows the complete statistical request}
\label{exec21:grfr-contract}
The initial proposal profile is one binary action and complete independent customer rows, with a fixed outcome horizon and all-population AIPW target. Clusters, sample weights, continuous actions, instrumental variables, survival, multiarm comparisons and repeated policies need separate profiles. The source exposes several of these methods, but their presence in a dispatcher does not transfer their assumptions or uncertainty guarantees to the initial bank use. A cohort chosen after observing treatment or outcome also changes the evaluation problem; the documented subset warning recommends selection using baseline features. Adaptive selection using fitted effects needs its own honest evaluation design.~\citep{exec21_grf_reuse}

Prediction outputs need equally careful typing. With no \texttt{newdata}, the causal-forest wrapper uses stored or recomputed out-of-bag predictions; explicitly supplying the training matrix as new data follows a different prediction path. Those calls are not interchangeable evidence of sample separation. The documented CATE variance, R-loss diagnostic and Monte Carlo excess error answer different questions. Increasing tree count can address forest simulation instability without creating exchangeability, recovering unsupported customers or resolving dependence. Model serialization must also account for the stored original observations and nuisance vectors.~\citep{exec21_grf_reuse}

The future worker should retain target identity, actual estimator, customer keys, column order, units, time horizon, warnings, source/dependency versions and nuisance provenance. It must distinguish unavailable execution, an unsupported request, a refused profile, an audited nonidentification finding and a failed computation. Failed or nonfinite output must not silently publish a partial estimate. Qualification would compare the original package against exact score and aggregation fixtures, test OOB and new-data routing, preserve warnings and refusal behavior, and check persistence and process failures. Public generated laws can then assess bias, interval coverage, widths and failed requests before an expensive bank trial. The small examples above establish the logic of that experiment, not its outcome.

Porting the wrapper or embedding its native core remains an unattractive initial route because it transfers nuisance, targeting, uncertainty and lifetime responsibilities without an observed benefit. Bayesian bcf, XBCF and flexBCF remain substantive comparators with different priors, uncertainty and restoration contracts, as developed in the preceding comparison. An original R worker is therefore the preferred prototype candidate, conditional on its environment and rights review. The measured comparison must include validation, auxiliary fits, training, scoring, serialization, human review and support costs for accepted outputs. No present evidence establishes a speed advantage, lower total cost or bank value for any route.~\citep{exec21_grf_reuse,exec21_cxx_bcf,exec21_cxx_xbcf,exec21_cxx_flexbcf}



% SOURCE docs/release/2026-10-29/source/docs/writing-repair/2026-10-29/src/6d9ce80d3a2938d8.tex CHAR 0
\section{Julia effect estimators and the meaning of a reusable result}
\label{exec21:effects-scope}

A servicing manager asks whether an additional contact is worth its cost for
each risk group. A numerical answer requires more than a treatment-effect
method name. The population average, the average among previously contacted
customers, a conditional contrast and a coefficient from residual regression
can differ even when every conditional model is known. Missing outcomes can
change the averaging population again. A reusable interface must carry this
meaning from the business question through the numerical calculation.

The selected Julia sources offer several distinct building blocks.
CausalELM~0.8.0 supplies g-computation and effect learners; CausalForest~0.1.5
implements tree-based contrasts; TMLE~0.20.5 supplies targeted and one-step
estimation; and CausalTables~1.3.4 represents causal data and interventions in
supplied structural models. Their selected source revisions precede the
7~October~2026 cutoff. Their MIT notices permit substantial reuse subject to
retaining the specified notices, but do not establish dependency clearance or
statistical validity. The following comparison concerns inspected source
paths, not native execution or a tested combination of these packages.
\cite{exec21_effects_causalelm,exec21_effects_causalforest,exec21_effects_tmle,exec21_effects_causaltables}

\subsection{An average effect is a population-specific quantity}
\label{exec21:effects-target}

Let $W$ denote pre-contact risk information, $A\in\{0,1\}$ the contact action,
and $Y(a)$ the monetary outcome under action $a$. Under consistency,
conditional exchangeability and treatment positivity, write
$e(w)=\Pr(A=1\mid W=w)$ and
$\tau(w)=E\{Y(1)-Y(0)\mid W=w\}$. The population average treatment effect is
$E\{\tau(W)\}$; the treated-population effect is
$E\{\tau(W)\mid A=1\}$. A prediction vector containing estimated
$\tau(W_i)$ is useful, but does not by itself provide uncertainty for a
selected subgroup or a policy chosen using that vector.

In the inspected CausalELM g-formula path, ATE, ITT and CATE labels use
counterfactual predictions at all supplied covariate rows; the other branch
uses treated rows. The g-computation wrapper averages the contrast, while the
S-learner retains a contrast vector. The T-learner fits each treatment arm and
subtracts its predictions on the supplied covariates. Thus an ITT label only
has an assignment-effect interpretation when the supplied treatment actually
represents assignment; it cannot convert receipt of contact into randomized
offer assignment. Categorical coding, the action clock and the averaging
population must be explicit inputs. \cite{exec21_effects_causalelm}

The same source's double-machine-learning path forms outcome and treatment
residuals and averages foldwise least-squares coefficients. Its fold
constructor slices contiguous rows. Neither the name nor contiguous slicing
supplies a dependence argument for repeated customer records. Even with
oracle nuisance functions and independent observations, a partially linear
coefficient need not be the population average under heterogeneous effects.
To see the distinction, take
$Y=m_0(W)+A\tau(W)+\varepsilon$ with
$E(\varepsilon\mid A,W)=0$. If $m(W)=E(Y\mid W)$, then
$Y-m(W)=(A-e(W))\tau(W)+\varepsilon$. Consequently the population residual
coefficient is
\begin{equation}
\label{exec21:effects-residual}
 \theta_{\rm res}
 =\frac{E[(A-e(W))(Y-m(W))]}{E[(A-e(W))^2]}
 =\frac{E[e(W)(1-e(W))\tau(W)]}{E[e(W)(1-e(W))]}.
\end{equation}
This is an overlap-weighted average in this model. A constant effect or
additional conditions can make it agree with the ATE; cross-fitting alone
does not do so. This derivation concerns the oracle population target and is
not a claim that a fitted implementation attains it in finite samples.
\cite{exec21_effects_causalelm,exec21_effects_reference}

Return to the stipulated servicing population used for the R forest in
Section~\ref{exec21:grfr-target}: the two risk groups have equal population
mass, effects $(-4,12)$ dollars and treatment propensities $(1/2,1/10)$.
Using the same law makes the target comparison across languages explicit;
it supplies no independent empirical corroboration. Exact enumeration gives
\begin{equation}
\label{exec21:effects-targetexample}
 \tau_{\rm ATE}=4,\qquad
 \tau_{\rm ATT}=-\frac43,\qquad
 \theta_{\rm res}=\frac4{17}.
\end{equation}
At an assumed cost of USD~2 per additional contact, the population margin is
USD~2. Subtracting that cost from the residual coefficient instead gives
$-30/17$ dollars. These quantities answer different questions. A proposed
product should refuse a request to relabel the coefficient as the
population-wide incremental value. A separately evaluated targeted policy
could be more useful, but selecting the profitable subgroup and estimating
its value require their own information and uncertainty contract.
\cite{exec21_effects_reference}

\subsection{Source-level uncertainty checks before adoption}
\label{exec21:effects-uncertainty}

The inspected CausalELM generic null routine samples treatment independently
and uniformly from its distinct observed values, or uniformly over its
observed range for a continuous treatment. This does not preserve the
observed treatment count. It also does not reproduce a nonuniform propensity
law or a cluster-randomized design without additional construction. For four
customers assigned exactly one contact, independent fair binary draws retain
that design with probability $4/16=1/4$ and create an empty treatment or
control arm with probability $2/16=1/8$. This exact calculation identifies an
assignment-law mismatch, rather than estimating a package error rate.
\cite{exec21_effects_causalelm,exec21_effects_reference}

There is also a direct dimensional problem in the selected source. For null
statistics $z_1,\ldots,z_B$ and reported statistic $t$, its scalar
\texttt{stderr} expression is
\begin{equation}
\label{exec21:effects-scale}
 s_{\rm code}=\frac{\sum_{b=1}^{B}(t-z_b)^2}{(B-1)\sqrt B}.
\end{equation}
If $t$ and $z_b$ are monetary effects, this expression has squared monetary
units. Multiplication of all outcomes by a positive constant $c$ multiplies
it by $c^2$, whereas a standard error must scale by $c$. With
$(z_1,z_2,z_3,z_4)=(-1,-1,1,1)$ and $t=2$, the expression is $10/3$.
Changing dollars to cents multiplies it by 10,000 rather than 100. This is a
source-level defect in the expression at the selected revision. It does not
require a Julia run to establish the algebra, and it does not establish the
finite-sample behavior of every other package function.
\cite{exec21_effects_causalelm,exec21_effects_reference}

Taking a square root in one location would not settle the larger inference
question. The sampled null law, the centering, the effect being tested and
the relationship between repeated randomization draws and sampling
uncertainty would still need justification. The inspected interval routine
searches a finite grid with tail comparisons against the same null draws;
that code alone is insufficient to claim a valid general confidence-interval
inversion. The proposed reuse decision is therefore to reject unqualified
reuse of this uncertainty path, while retaining the point-estimation family
as a candidate for a separately justified implementation.
\cite{exec21_effects_causalelm}

CausalForest illustrates a different boundary. Its interface exposes
bootstrap or subsampling choices and separate split and leaf samples under
the honest-tree option. However, its outcome-centering forest is fitted
before these per-tree splits. The inspected prediction paths form
within-leaf treated and control weights and return centered-outcome
contrasts. They return point predictions, rather than a calibrated covariance
object. Honesty of the tree split is consequently not sufficient evidence
that all nuisance fitting is independent of evaluation, nor evidence of
equivalence to the separately reviewed \texttt{grf} implementation. A native
comparison should examine the actual nuisance training rows, treatment
coding, arm support in leaves, random-number state and failure behavior
before considering any shared inference contract.
\cite{exec21_effects_causalforest}

\subsection{Targeted estimation, missingness and influence scales}
\label{exec21:effects-tmle}

TMLE exposes nuisance-model choices, outer resampling and propensity
truncation. The selected noncollaborative default uses no outer resampling;
a requested cross-validation design must be supplied and preserved. The
implementation's \texttt{weighted} constructor default is true although
adjacent prose says false, another reason to record actual resolved
configuration. Truncation prevents numerical division by a very small
estimated probability but cannot create absent treatment support. For an
identified ATE in an independent sampling model, the familiar influence
expression clarifies what an adapter must carry:
\begin{equation}
\label{exec21:effects-influence}
\begin{aligned}
 D_i={}&\frac{A_i}{e(W_i)}\{Y_i-m_1(W_i)\}
       -\frac{1-A_i}{1-e(W_i)}\{Y_i-m_0(W_i)\}\\
       &+m_1(W_i)-m_0(W_i)-\psi .
\end{aligned}
\end{equation}
The relevant source combines outcome-residual and counterfactual-mean
components. It stores the standard deviation of the estimated influence
values together with their number, and passes both to the testing
interface. This scale is not already the standard error of the effect.
Conditional asymptotic inference additionally needs suitable nuisance
remainder control, positivity, moments and a sampling-law argument; the
presence of an influence vector does not prove those conditions.
\cite{exec21_effects_tmle}

For example, influence values $(-3,-1,1,3)$ have sample variance $20/3$.
Under the independent mean calculation, the variance of the average is
$5/3$, since $n=4$. Exporting the influence standard deviation as an effect
standard error would miss the factor $\sqrt n$. Conversely, dividing a
quantity that is already a standard error by $\sqrt n$ again would understate
uncertainty. For joint outputs, the source stacks influence vectors and
computes their covariance. A bridge must preserve the identity and order
of contributing customers; equal vector length is insufficient when
different outcomes select different complete cases. Cluster or calendar
dependence requires an appropriate covariance construction, not a new field
name on the independent calculation. \cite{exec21_effects_tmle,exec21_effects_reference}

The missingness documentation is explicit: rows missing variables needed
for estimation are dropped. Without cross-validation, propensity fitting
may use additional records with observed treatment and confounders even
when outcomes are missing. The source's final fluctuation dataset nevertheless
uses complete relevant records. This is useful data reuse, but is not a
general missing-outcome correction. \cite{exec21_effects_tmle}

Let $R$ indicate observation of the outcome. Consider independent customers
with equally frequent risk groups, randomized contact probability $1/2$,
effects $(-4,12)$ and outcome-response probabilities $(9/10,1/10)$ by risk
group. Suppose response is independent of potential outcomes and contact
conditional on risk, so missing at random holds in this example and response
positivity holds. The complete cases have group probabilities $(9/10,1/10)$,
giving
\begin{equation}
\label{exec21:effects-response}
 E\{\tau(W)\}=4,\qquad
 E\{\tau(W)\mid R=1\}=-\frac{12}{5}.
\end{equation}
Correct conditional response weighting recovers 4 in the exact known law.
At cost USD~2, the population margin is USD~2 and the complete-case margin
is USD~$-22/5$. The example does not claim these are observed package
outputs. It shows why improving a treatment-propensity fit cannot, by
itself, repair a changed marginal population. A missing-data-aware
identification and estimation design is required for the original request.
\cite{exec21_effects_reference}

Case-control sampling is also a distinct design. The inspected TMLE
interface accepts a population prevalence as a scalar and constructs
prevalence weights. Uncertainty in an externally estimated prevalence is
not an input to those inspected interfaces. A banking adapter must therefore
state whether prevalence is assumed fixed or estimated and how its
uncertainty enters a proposed joint analysis. It should not silently reuse
case-control weights for nonresponse, population transport or clustered
sampling, whose information and identifying assumptions differ.
\cite{exec21_effects_tmle}

\subsection{Supplied causal models are valuable public references}
\label{exec21:effects-scm}

CausalTables supplies a useful route to inexpensive comparisons before bank
experiments. Its intervention functions modify treatment, recompute
summaries and draw counterfactual responses from a supplied structural
causal model. Its \texttt{cfmean} and \texttt{cfdiff} functions approximate
intervention means by Monte Carlo. Despite efficiency-bound wording in
their documentation, the inspected return tuples contain only the mean.
A reference-generation service must separately retain the generating
equations, intervention, random seed, Monte Carlo error and dependence
structure. A simulated network summary does not establish a variance formula
for observations from a bank network. \cite{exec21_effects_causaltables}

The distinction between supplied truth and observational identification is
not semantic. Let $U$ be a fair binary variable and consider two structural
models. In both, $A=U$. In the first, $Y=A$; in the second, $Y=U$. Both
produce the same observed law, with probability $1/2$ on each of
$(A,Y)=(0,0)$ and $(1,1)$. Yet the intervention contrast is 1 in the first
model and 0 in the second. A simulator can reproduce either contrast once
its model is supplied. The observed distribution alone does not choose
between them. At an illustrative cost of $1/2$, the corresponding margins
are $1/2$ and $-1/2$. \cite{exec21_effects_reference}

TMLE's CausalTables extension converts treatment columns to categorical
arrays and identifies an estimand using the causes recorded in the supplied
table. This is an explicit and promising connection, provided that level
order, graph semantics and the provenance of the causal declarations survive
conversion. It does not verify that the declared graph is true. The
matching CausalInference revision supplies separately reviewed graphical
interfaces; that evidence should be reused for its established scope rather
than treated as qualification of heterogeneous-effect estimation.
\cite{exec21_effects_tmle,exec19_otherid_julia}

\subsection{A bounded adoption path and the wider language search}
\label{exec21:effects-integration}

The proposed architecture is an isolated Julia worker with a pinned
environment and an explicit estimand/result schema. TMLE is a candidate for
targeted mean estimation and CausalTables for supplied-model reference
generation. A first public comparison should use exact finite distributions
to verify scalar targets, treatment-level order, outcome-unit conversion,
complete-case selection and graph-to-estimand conversion. Subsequent native
trials should compare influence-vector scaling, fold membership, fitted
nuisance behavior and refusal on absent support. CausalELM's selected
uncertainty path needs a justified repair before inclusion; CausalForest's
point predictions need their own inference qualification. Porting either
path before understanding these contracts would copy its unresolved
scientific meaning into a new language.

The library, CLI and MCP representations should carry the action/assignment
distinction, target population, contributing unit IDs, missingness rule,
fitted-model configuration, uncertainty kind and all unavailable results.
An agent can propose mappings and inspect contradictory declarations, but
must leave unsupported identification or uncertainty requests unresolved.
This gives the business user a meaningful answer when a requested population
effect cannot be supplied, and permits genuine comparison of methods that
do answer the same question.

Rust coverage remains incomplete. Separate official repository and crate
queries identify leads including \texttt{perpetual}, \texttt{causalkit},
\texttt{oxicuda-causal} and \texttt{propensity-score}, alongside graph and
event-ordering projects. These are metadata leads; current descriptions are
not proof of functionality at the cutoff. The queries inspect only their
first pages, and an exact treatment-effect phrase query returning zero does
not establish absence. Selected pre-cutoff source revisions, licenses,
estimators and inference paths must be inspected before any capability or
reuse conclusion. Wider language coverage and package-paper originals also
remain to be established. \cite{exec21_effects_discovery}

These public comparisons can determine which Julia integrations merit a
native prototype and which assumptions a later banking experiment must test.


% SOURCE docs/release/2026-10-29/source/docs/writing-repair/2026-10-29/src/094e4fdf55f43b37.tex CHAR 0
\section{A restricted CausalML route for heterogeneous offer effects}
\label{exec14:causalml}

Flexible response prediction becomes useful to a bank only when its contrast
matches the offer, population and information available at the decision date.
The inspected Uber CausalML revision declares version 0.17.0, Python at least
3.11, eighteen required dependencies and an Apache-2.0 root notice. Its
T-learner and R-learner provide plausible reusable components for a static
offer comparison. This is a source-based selection, without an installed
environment or measured interval performance \citep{exec14_causalml_primary}.

The T-learner fits separate response models for each treatment and the shared
control. For a binary offer $A$, write
$m_a(x)=E(Y\mid A=a,X=x)$. Under consistency, conditional exchangeability,
overlap and a suitable sampling law, the conditional mean contrast is
$\tau(x)=m_1(x)-m_0(x)$. It is a contrast of conditional means, not an identified
individual customer's two potential outcomes. The target population then
determines how these contrasts should be averaged.

An exact stipulated account population makes the difference visible. Two
pre-offer strata have equal population mass and offer probabilities $1/5$
and $4/5$. Their mean recoveries without an offer are USD~100 and USD~120;
with an offer, USD~105 and USD~135. Suppose exchangeability holds within each
stratum. The conditional effects are 5 and 15, hence
\[
\operatorname{ATE}=\tfrac12(5)+\tfrac12(15)=10,\qquad
\operatorname{ATT}=\tfrac15(5)+\tfrac45(15)=13,\qquad
\operatorname{ATC}=\tfrac45(5)+\tfrac15(15)=7.
\]
The observed treated mean is 129 and the untreated mean is 104, giving an
unadjusted difference of 25. All four numbers are correct for their stated
calculations, but only one is the population ATE. The adapter must retain
the population measure instead of treating the presence of a number named
``effect'' as sufficient \citep{exec14_causalml_oracle}.

If the offer costs USD~8 and the imposed means are known, offering only to the
second stratum gives expected gain $\tfrac12(15-8)=3.5$ dollars per eligible
account, compared with 2 dollars for offering to everyone. This is a planning
calculation with known means. Selecting a rule from fitted effects introduces
uncertainty and selection; these dollar amounts do not establish a learned
policy's value or observed bank ROI.

The inspected implementation returns contrast columns in its treatment-label
order. The wrapper should use that returned mapping, preserve account and
feature IDs, and state treatment-minus-control units. The initial route should
remain binary: in the inspected multi-treatment path, the ATE point average
uses prediction rows while the variance terms use each treatment/control
subset. The populations therefore require a separate analysis before a
multi-arm uncertainty contract is promised \citep{exec14_causalml_primary}.

Uncertainty also has several meanings. The selected CATE interval path uses
prediction percentiles from refitted samples; the default ATE expression
combines arm residual variances and fitted-effect variation. The base bootstrap
resamples rows. None of these operations automatically accounts for household
clusters, repeated account histories or common calendar shocks. With a treated
fraction of $1/100$ and a resample of size 100, the probability of selecting
no treated row is $(99/100)^{100}$, about 0.366. A failed arm cannot be silently
discarded to improve reported interval availability. Changing to stratified
or block resampling changes the procedure and needs its own declared target
and analysis \citep{exec14_causalml_primary,exec14_causalml_oracle}.

The R-learner offers a related computational alternative. Its inspected caller
forms outcome residuals and a weighted regression with pseudo-outcome
$(Y-\widehat m(X))/(A-\widehat e(X))$ and weight
$(A-\widehat e(X))^2$. In the deterministic conditional-mean example above,
substituting the true $m$ and $e$ makes that pseudo-outcome equal to the stratum
effect for both treatment values. The identity explains the intended
orthogonalization; it does not validate learned nuisance functions, delegated
propensity folds or dependence-aware cross-fitting. Those delegated paths
must be reviewed before a broader route is enabled.

Observed uplift curves are another distinct object. The inspected metric
uses true effects if supplied and otherwise uses treated/control observed
means with interpolation. A curve can be computable under confounding and
still fail to identify policy value. An ordinary exact stratified estimator
therefore remains a useful comparator, alongside separately reviewed EconML
and grf paths.

The preferred proposal is an optional pinned T-learner service with explicit
base regressor, binary arms and independent whole accounts. Preserve every
seed, environment identity, population weight, pre-offer feature clock and
outcome horizon. Reject missing arms, unsupported sampling, stale row mappings,
nonfinite output and unsupported targets with a reason. Keep time-varying
treatment, latent-state inference and multi-arm uncertainty outside this first
contract. Full adoption would import several scientific contracts at once;
porting would assume their maintenance without a demonstrated need. Executable
admission remains deferred until input/import paths, dependency notices and
target, column, interval and failure fixtures have been reviewed and run.
\label{exec14:causalml-end}


% SOURCE docs/release/2026-10-29/source/docs/writing-repair/2026-10-29/src/bd73ca27d872f211.tex CHAR 0
\section{Choosing a pinned statsmodels reference}
\label{sec:exec24-statsmodels-reuse}
The earlier Python comparison explains why a bank must distinguish ATE, ATT and ATU, preserve the population in each moment, expose propensity clipping and follow uncertainty through the final calculation. These requirements favor an interpretable parametric reference when its assumptions match. They do not give every model in a broad statistical library an accepted causal interpretation. Retain the selected statsmodels treatment-effect and GMM paths as references, with a pinned original Python worker as a conditional prototype route. A whole-library port has no demonstrated benefit here.~\citep{exec20_python_effect_source,exec20_python_effect_reference}

The identity is commit \texttt{8278e2d218cc85bac2c7af02feb9a19a0e499b04} of 6 October 2026. Its project metadata declares a dynamic version; an installed package cannot supply an unverified release identity for this source. The authenticated build uses \texttt{mesonpy}, Meson-Python, Cython at least 3.0.13 below 4, NumPy at least 2.0 below 3, SciPy at least 1.13 below 2, and setuptools-scm at least 9.2 below 10. Those build constraints differ from runtime constraints.~\citep{exec24_statsmodels_distribution}

Runtime metadata permits NumPy from 1.23.5 below 3, SciPy from 1.8 except 1.9.2, pandas from 1.4 except 2.1.0, patsy from 0.5.6, packaging from 21.3 and formulaic from 1.1.0, with Python at least 3.10. Test, documentation, development and Polars extras are separate. These ranges are not a solved environment. The complete pinned tree contains project metadata and a Meson build file but no root \texttt{setup.py}; that filename's failed request does not establish package unavailability.~\citep{exec24_statsmodels_distribution}

Direct BSD terms require applicable notices, conditions and disclaimers in source or binary redistribution and restrict endorsement through project or contributor names. They do not clear all dependencies, data or distribution arrangements. A maintainer must bind the actual artifacts, optional extras and license inventory. A commit date records source custody, not a maintenance or support guarantee.~\citep{exec20_python_effect_source,exec24_statsmodels_distribution}

\subsection{Follow uncertainty to the final result}
The selected class offers IPW, regression adjustment, augmented IPW, weighted-outcome augmented IPW and IPW regression adjustment with linear OLS outcome models. Target population affects weighting and averaging. A preliminary outcome fit's robust covariance label does not determine final joint GMM uncertainty. Section~\ref{exec20:pyfx-covariance} follows the selected final moment path and stored covariance settings: the wrappers retain the default moment-covariance route, and a preliminary robust label does not expose the separate HAC configuration. Native qualification must test that complete path; clustered or HAC inference cannot be inferred from a preliminary fit option.~\citep{exec20_python_effect_source}

Prior exact examples remain behavioral oracles. Their stipulated population gives ATE 50, ATT 65 and ATU 35, so relabeling changes a decision. Clipping a propensity can change an augmented score under an incorrect outcome model; it does not create overlap. The covariance example distinguishes scalar centering from per-column centering when moment means are nonzero, while exactly solved zero means remove that difference. These authored examples are not newly observed native defects.~\citep{exec20_python_effect_reference}

Freeze target, sample keys, treatment coding, nuisance predictions, clipping and dependence assumptions before native qualification. Compare final moments, Jacobian, residuals, rank and covariance against an independent small reference. Test row permutations and grouped control/treated order; preserve singularity, nonconvergence, warnings and unavailable uncertainty. Point availability and justified uncertainty are separate outputs. An unsupported common branch or calendar law cannot be repaired merely by selecting a familiar robust label.~\citep{exec20_python_effect_source,exec20_python_effect_reference}

\subsection{Decision, costs and accountable roles}
Preserve the original Python engine behind a typed boundary as the preferred conditional prototype. A whole-library R, Java or C++ port duplicates formula processing, numerical behavior, uncertainty and maintenance. Explicit estimating equations or an appropriate DoubleML profile remain comparators according to target and nuisance assumptions. The method owner must approve the sampling law, the integration maintainer must own conversion and upgrades, a rights reviewer must resolve distribution, and an evaluation owner must freeze tests and costs. These are obligations, not external acceptances already obtained.~\citep{exec20_python_effect_source,exec24_statsmodels_distribution}

Include source and rights review, dependency locking, data and clock conversion, nuisance fitting, optimization, covariance/resampling, failures, retries, storage and support in costs. No supplied price or return follows from this source review. The proposal reference choice can close with \texttt{integration\_allowed} false. Revisit native use only with resolved environment and rights, registered semantic/failure tests and matched complete costs. Native performance, bank benefit and general coverage remain open.~\citep{exec20_python_effect_source,exec24_statsmodels_distribution}


% SOURCE docs/release/2026-10-29/source/docs/writing-repair/2026-10-29/src/471448bf7ebe65e7.tex CHAR 0
\section{R tmle: preserve the target and audit uncertainty before reuse}
\label{sec:exec24-tmle-reuse}

A banking team considering a single repayment offer may need the change in
repayment probability, the ratio of repayment risks, or the effect among
customers who actually received the offer. These targets can lead to different
allocation decisions. The selected R \texttt{tmle} 2.1.1 distribution supplies
binary point-treatment machinery, outcome-missingness adjustment, ATT and ATC
paths, and a binary intermediate-variable route. It is a useful reference for
an eventual native R worker. It does not, merely by accepting repeated customer
identifiers, establish identification for a sequence of adaptive interventions.
Our earlier inspection of missing treatment and baseline covariates remains
applicable: these are rejected, and the no-variation remapping is a mean-outcome
route rather than reconstruction of unknown assignments.
\citep{exec20_rfx_tmle}

The selected CRAN mirror revision is dated 30 August 2025 and precedes the
literature cutoff. We inspect its actual treatment function, help text,
dependency declaration, rights file and citation metadata. The latter identifies
a 2012 package article, but is not a substitute for reading its original.
The comparison here rests on the pinned technical implementation and authored
arithmetic; it does not borrow uninspected results from the package bibliography.
\citep{exec24_tmle_sources}

\subsection{Point-treatment meaning, weighting and the population clock}
Let $W$ contain information available before the offer, $A\in\{0,1\}$ its
assignment, $Y$ the later repayment indicator, and $\Delta$ its observation
indicator. Under consistency, conditional exchangeability, treatment positivity
and an appropriate outcome-observation assumption, an identified point-treatment
mean may be represented by an average of $Q_a(W)$, where $Q_a$ is the conditional
outcome regression at treatment $a$. These are substantive premises for the
proposal, not conclusions of a successful regression. Sampling weights $w_i$
produce the empirical plug-in targets
\[
 \mu_a=\frac{\sum_i w_i Q_a(W_i)}{\sum_i w_i},\qquad
 \mathrm{RD}=\mu_1-\mu_0,\quad
 \mathrm{RR}=\frac{\mu_1}{\mu_0},\quad
 \mathrm{OR}=\frac{\mu_1/(1-\mu_1)}{\mu_0/(1-\mu_0)}.
\]
The source normalizes observation weights to sum to the number of rows and
uses corresponding weighted means. Its parameter helper calculates these
contrasts and distinct influence-function expressions. The adapter must report
the actual target, denominator, observation rule and weight interpretation.
\citep{exec24_tmle_sources}

In an authored two-profile example, weights $(1,3)$ and regression predictions
$Q_0=(0.1,0.4)$, $Q_1=(0.2,0.5)$ give $\mu_0=13/40$, $\mu_1=17/40$,
$\mathrm{RD}=1/10$, $\mathrm{RR}=17/13$ and $\mathrm{OR}=459/299$.
At a stipulated 200 USD per incremental repayment, the difference gives
20 USD per offer before costs; the ratio itself is not a dollar benefit.
Neither a fitted conditional mean nor an externally supplied prediction matrix
establishes the causal assumptions. A later prototype must preserve training,
fold, prediction and customer identities and qualify the actual nuisance
estimation and uncertainty route.
\citep{exec24_tmle_contract}

Observation and customer populations also differ. One customer with one row
and effect zero, and another with three rows each having effect ten, give a
row-average effect of $7.5$ and an equal-customer average of $5$. The inspected
helper averages influence-function contributions within identifiers and divides
the variance of those means by the number of identifiers. Whether that aligns
with a weighted row target requires its sampling law and cluster-size argument;
an identifier field alone is insufficient.
For row scores $(-3,1,1,1)$ in clusters of sizes one and three, the cluster means
are $(-3,1)$ and their sample variance divided by two is $4$. For comparison,
the cluster-sum contributions scaled by the observed cluster-to-row ratio are
$(-1.5,1.5)$, whose sample variance divided by two is $9/4$. This is exact
arithmetic illustrating different contributions, not a proof of native
miscalibration or a universal prescription for dependent observations.
\citep{exec24_tmle_sources,exec24_tmle_contract}

\subsection{Intermediate-variable targets and a specific interval discrepancy}
When the supplied intermediate variable has more than one observed value, the
source constructs intervention predictions at its two binary levels and uses
joint treatment, intermediate-variable and outcome-observation probabilities.
A controlled intervention on a mediator needs a defensible action and the
corresponding sequential exchangeability and support assumptions. Conditioning
on an observed mediator value is not automatically the same intervention.
Truncating a probability of $0.01$ to $0.1$ changes an inverse contribution from
100 to 10; it does not create unobserved interventions or prove identification.
The source's explicit bounds and retained rows belong in the result contract.
\citep{exec24_tmle_sources,exec24_tmle_contract}

A concrete static discrepancy matters for this route. Its bootstrap matrix is
named, in order, \texttt{EY0}, \texttt{EY1}, \texttt{ATE}, \texttt{logRR} and
\texttt{logOR}. The source fills the fourth column with log risk ratios and
uses column four for the risk-ratio variance, but constructs its risk-ratio
bootstrap interval by exponentiating quantiles of column three. The ordinary
point-treatment branch uses column four for this interval. We verified the
column definitions, assignments, branch condition and output statements; this
is a mismatch in the selected source, not an observed package execution.
\citep{exec24_tmle_sources}

The difference is visible without running that implementation. If every row of
an authored bootstrap fixture has risks $0.1$ and $0.2$, its risk difference is
$0.1$ and its risk ratio is $2$. Exponentiating the constant difference gives
$\exp(0.1)\simeq1.10517$ instead of $2$. Thus column substitution changes the
scientific quantity even when numerical routines return finite values. The
proposed integration refuses the affected interval pending a separately
reviewed correction or qualified successor. A repair must use named quantities,
check both intermediate-variable levels and all interval fields, and preserve
the unmodified original with the evidence. We do not patch upstream code or
silently present our arithmetic as its native output.
\citep{exec24_tmle_contract}

\subsection{Bootstrap counts, missing outputs and uncertainty admission}
Both inspected branches iterate from one through $B$. Iterations below $B$
resample identifiers, while the last uses the original sample so returned fits
have the original identities. Quantiles are calculated over the full matrix.
Consequently the inspected loop contains $B-1$ resampled rows and one original
row, which must be distinguished from a request for $B$ independent resamples.
For an authored vector $(0,1,100)$, where the last value represents the original
sample, linear-interpolated 2.5 and 97.5 percent quantiles are $0.05$ and $95.05$.
Using only the two stipulated resamples gives $0.025$ and $0.975$. These numbers
isolate a counting convention; they do not measure empirical coverage.
\citep{exec24_tmle_sources,exec24_tmle_contract}

The inspected variance call does not remove missing values, while its quantile
calls explicitly do. With the illustrative vector $(1,\mathrm{NA},3)$ the
variance is missing but the finite-only quantile interval is $[1.05,2.95]$.
A successful interval field cannot hide its incomplete denominator. The future
worker must return requested, attempted, original-sample, finite and failed
counts, and identify the actual nuisance-fitting and targeting steps repeated.
The intermediate-variable bootstrap samples identifiers separately from the two
observed mediator groups and then includes all rows for each sampled identifier.
Identifiers spanning both groups require explicit scrutiny; disjoint resampling
populations cannot be presumed from a row-level mediator label.
\citep{exec24_tmle_sources,exec24_tmle_contract}

The proposal therefore separates influence-function intervals, targeted
bootstrap intervals, sensitivity scenarios and identification failures. Boundary
risks, zero ratio denominators, unsupported missingness, nuisance failures,
truncation, nonconvergence and incomplete replicates need typed outcomes.
Neither an interval label nor the documentation's suggested high iteration
count proves coverage for arbitrary bank panels or adaptive policies.
\citep{exec24_tmle_contract}

\subsection{Rights, alternatives and a bounded adoption decision}
The DESCRIPTION declares \texttt{BSD\_3\_clause + file LICENSE | GPL-2}; the
selected LICENSE file contains year and copyright-holder fields. The treatment
source footer separately describes educational, research and not-for-profit
permission and directs commercial inquiries elsewhere. We preserve these
texts as a rights question for the selected distribution, without deciding their
legal precedence. No commercial permission is inferred from a package index.
Declared dependencies include \texttt{glmnet} and \texttt{SuperLearner}, with
optional learners and diagnostics; an installed dependency closure, its rights,
platform behavior and maintenance commitments remain to be qualified.
\citep{exec24_tmle_sources}

Our current disposition is reference. Conditional integration through the
original R worker is preferable to assuming that a whole-library port preserves
fitted objects, weighting and targeting semantics. A narrow independent
finite-population or analytical oracle can establish target and refusal
expectations. A separately qualified longitudinal estimator answers a different
request; it is not an alias for this point-treatment interface. Whole-library
porting is deferred until demonstrated demand and costs justify it. The
scientific owner must approve target, sampling and missingness assumptions; the
adapter owner must preserve identities and typed failure; the rights/platform
owner must resolve distribution and environment; and the evaluation owner must
preregister public positive, adverse, boundary and fault fixtures.
\citep{exec24_tmle_contract}

Illustrative total cost of 2,400 USD over 120 requests, 100 conformant results
and 80 useful results is respectively 20, 24 and 30 USD per denominator. These
are authored inputs, not measured prices or a business case. Rights review,
engineering, dependencies, compute including failures, storage, recovery and
maintenance belong in the same cost record. Public or generated experiments
can justify the next investment before expensive bank data work. They cannot
supply unknown assignment truth, universal hidden-confounding identification
or arbitrary dependent-data coverage. The interval discrepancy has a concrete
repair and validation route; scientific identification and commercial rights
remain explicit conditions, not promises of automatic closure.
\citep{exec24_tmle_contract}


% SOURCE docs/release/2026-10-29/source/docs/writing-repair/2026-10-29/src/d47fc913470c809b.tex CHAR 0
\section{Delicatessen reuse: estimating equations, numerical roots and the sampling unit}
\label{sec:exec24-delicatessen-reuse}

A bank comparing servicing actions needs more than a fitted effect and a
standard error. It needs to know which equations define that effect, which
observations supply independent information, and which failures prevent an
uncertainty statement. Delicatessen is a useful reference for making these
choices explicit through stacked estimating equations. The proposed use is
a bounded numerical reference and possible future isolated worker, not a
general certificate of causal identification.~\citep{exec24_delicatessen_source,exec24_delicatessen_choice}

Seven selected originals authenticate to the repository revision dated
11 August 2026. The version file reports 4.3. The inspected setup text
declares NumPy at least 1.18.5 and SciPy at least 1.9.0, and lists Python
3.8 through 3.14 classifiers; classifiers are not a tested environment
matrix or an enforced interpreter constraint. The top-level MIT text is
preserved. Installed-wheel equality, transitive rights, numerical backends
and supported environments remain separate qualification questions.
The source is inspected as text, without executing downloaded code.
\citep{exec24_delicatessen_source}

\subsection{The equation is the statistical contract; a root is a numerical result}
Let $\theta\in\mathbb R^p$ be the declared parameter and let
$\psi_i(\theta)\in\mathbb R^p$ be the contribution from independent unit
$i$. The estimating equation and its empirical derivative are
\[
 \sum_{i=1}^n\psi_i(\widehat\theta)=0,\qquad
 A_n=-\frac1n\sum_{i=1}^n\partial_\theta\psi_i(\widehat\theta).
\]
With $M_n=n^{-1}\sum_i\psi_i\psi_i^\top$, the conventional regular,
nonsingular calculation is
\[
 \widehat{\operatorname{Var}}(\widehat\theta)
   =\frac1n A_n^{-1}M_nA_n^{-\top}.
\]
This expression requires a justified expansion, a suitable limit law and
an identified target. The formula itself supplies none of those facts.
The selected estimator obtains the unit count from the contribution
array, computes bread and meat scaled by that count, and then divides
the assembled asymptotic covariance by the count once more.
\citep{exec24_delicatessen_source}

The default solver is Levenberg--Marquardt, with 5,000 iterations and
tolerance $10^{-9}$. The selected built-in root branches inspect their
solver's success state. A custom callable instead receives the equations
and initial values, and must return a non-null result; the estimator's
declared iteration and tolerance options are not passed to that callable.
A proposed worker must therefore record the actual solver contract and
check a scaled full-equation residual, feasibility and sensitivity to
starting values after the solver returns. A success flag is not an
independent proof that the requested equation has been solved.
\citep{exec24_delicatessen_source,exec24_delicatessen_choice}

For an authored mean example with observations $(0,0,2,2)$,
$\psi_i(\theta)=Y_i-\theta$ has root $\widehat\theta=1$.
A custom routine returning $0$ has average residual $1$, even though
it has returned a numerical value. The proposed residual test rejects
that candidate before constructing an interval. It must use declared
equation scales because a residual tolerance is otherwise changed by
arbitrary multiplication of one equation.~\citep{exec24_delicatessen_choice}

The subset option solves only selected coordinates, holding the others
at their supplied initial values. The covariance calculation still uses
the full estimating-equation system. This is appropriate only when the
held coordinates have already been solved consistently with the declared
stack, or a separate derivation justifies a different treatment. It is
not a switch that turns estimated nuisance parameters into known constants.
The proposed result retains all residuals and the provenance of every
held nuisance value. Arbitrary externally trained learners and cross-fitting
require additional rate and sampling arguments beyond this selected API.
\citep{exec24_delicatessen_source,exec24_delicatessen_choice}

\subsection{A pseudoinverse does not identify an unconstrained direction}
The selected covariance builder defaults to a Moore--Penrose pseudoinverse;
disabling that option requests an ordinary inverse. A NaN bread matrix
causes the builder to return no covariance. Finite output from the
pseudoinverse branch must not be relabelled as scientific identification.
\citep{exec24_delicatessen_source}

Consider the authored system
\[
 \psi_i(\theta)=\begin{pmatrix}Y_i-\theta_1\\0\end{pmatrix},\qquad
 A_n=\begin{pmatrix}1&0\\0&0\end{pmatrix}.
\]
Every value of $\theta_2$ solves the second equation. For the four
observations above, the algebraic pseudoinverse sandwich has diagonal
$(1/4,0)$. The zero in the second position represents projection onto an
unconstrained direction, not exact knowledge of $\theta_2$. A scientific
adapter must distinguish an identified lower-dimensional target, an
identified set and an unidentified request. The latter cannot be repaired
by replacing a singular inverse with a finite matrix.
\citep{exec24_delicatessen_choice}

Numerical differentiation also requires its own contract. The selected
code offers several finite-difference choices and a separate automatic
differentiation route. The default increment is $10^{-9}$; adequacy
depends on parameter units, smoothness and numerical precision. A native
qualification should compare derivatives against analytic examples and
rescaled parameters, report rank and condition diagnostics, and retain
nonfinite states. Inspecting the dispatch does not qualify all derivative
implementations or nonsmooth user equations.~\citep{exec24_delicatessen_source,exec24_delicatessen_choice}

\subsection{Individual rows and independent clusters give different covariance}
The selected meat calculation multiplies the contribution array by its
transpose. It has no implicit knowledge that multiple account rows share
a branch shock or that adjacent months are dependent. Independence must
refer to the columns supplied to the numerical routine, with a scientific
justification for any aggregation or dependence correction.
\citep{exec24_delicatessen_source}

For the same four observations, individual mean scores are
$(-1,-1,1,1)$, the raw bread is $4$ and the raw meat is $4$, yielding
variance $4/4^2=1/4$. If the first pair forms one independent cluster
and the second pair another, the cluster-summed scores are $(-2,2)$.
Their raw meat is $8$, while the derivative of the full summed equation
remains $4$, yielding $8/4^2=1/2$. This is an arithmetic witness of the
missing cross terms, not evidence that two clusters suffice for reliable
asymptotic coverage.~\citep{exec24_delicatessen_choice}

Supplying cluster sums as columns preserves the original full estimating
equation. Unequal cluster sizes require particular care: replacing sums
with unweighted cluster means can change the root and its target.
Cross-cluster correlation or serial persistence needs a separately
justified law; aggregation alone does not solve arbitrary dependence.
The inspected finite correction offers HC1, multiplying the meat by
$n/(n-p)$ and rejecting $n\leq p$. Its count is the supplied unit count.
This correction is not a general small-cluster or HAC justification.
\citep{exec24_delicatessen_source,exec24_delicatessen_choice}

\subsection{Standalone covariance needs an explicit shape and normalization contract}
The standalone covariance function documents an asymptotic covariance
that the caller divides by the number of observations. Its selected
one-parameter branch takes the first dimension of the contribution array
as that number. Consequently a vector of length $n$ and a one-row array
of shape $1\times n$ can take different normalization paths, although
both represent one equation with $n$ contributions. The main estimator
uses a transposed array to obtain its observation count.
\citep{exec24_delicatessen_source}

Following the selected formulas symbolically for the four-value example,
a vector supplies count $4$: normalized bread and meat are both $1$,
so standalone asymptotic variance is $1$ and the documented final division
gives $1/4$. A one-row array instead supplies count $1$: bread and meat
are both $4$, so the standalone return is $1/4$. Dividing that return by
the actual four observations gives $1/16$. The main estimator's one-row
route gives $1/4$. Under HC1, the standalone row branch also reaches
$n=p=1$ and refuses the correction. These are conditional static
control-flow deductions and independently authored arithmetic, not
observed native runs.~\citep{exec24_delicatessen_source,exec24_delicatessen_choice}

The proposed worker should use a canonical contribution shape, compare
both public routes against a symbolic oracle, and label asymptotic versus
estimator covariance explicitly. It should qualify one-parameter vectors,
one-row arrays, multi-parameter arrays, clusters and finite corrections
before relying on either route. Silently dividing every returned matrix
by a row count would conceal the disagreement.
\citep{exec24_delicatessen_choice}

\subsection{AIPW clipping changes the protection offered by the nuisance models}
For binary action $A$, outcome $Y$, propensity $e(W)$ and outcome
functions $m_a(W)$, the selected AIPW contribution for the treated mean
is algebraically
\[
 m_1(W)+\frac{A}{\widetilde e(W)}\{Y-m_1(W)\}-\mu_1,
\]
with the analogous control expression using $1-A$ and
$1-\widetilde e(W)$. The selected routine stacks the mean difference,
both action means, propensity equations and outcome equations. It chooses
a logistic outcome model for all-binary observed outcomes unless the
continuous option is set. The counterfactual design matrices are supplied
by the caller; equal dimensions do not establish correct action or
interaction coding.~\citep{exec24_delicatessen_source}

The clipping branch checks endpoint order and applies the clipping
operation. That local branch does not enforce strict endpoints inside
$(0,1)$ or establish positivity. A worker needs explicit domain and
finite-value checks and must record the truncation policy. Conditional
on $W$, and under consistency, exchangeability and positivity, the mean
of the treated pseudo-outcome is
\[
 m_1(W)+\frac{e(W)}{\widetilde e(W)}
       \{q_1(W)-m_1(W)\},
 \qquad q_1(W)=E[Y\mid A=1,W].
\]
Thus a correct outcome function still removes its residual term. But
clipping a correct propensity generally removes the protection of that
propensity when the outcome function is wrong. Clipping is not, by itself,
a definition of an overlap-population causal target.
\citep{exec24_delicatessen_source,exec24_delicatessen_choice}

An authored one-stratum population has $e=1/10$, action means
$q_1=4/5$ and $q_0=1/5$, with independent assignment to well-defined
potential outcomes. Deliberately omit action from an intercept-only
outcome model. Its population fitted mean is
$m_1=m_0=E[Y]=13/50$, so the nuisance score is solved even though the
conditional outcome model is wrong. Without clipping, the population
AIPW means recover $4/5$ and $1/5$. Clipping the propensity to $1/5$
instead gives
\[
 \widetilde\mu_1=\frac{53}{100},\qquad
 \widetilde\mu_0=\frac{77}{400},\qquad
 \widetilde\tau=\frac{27}{80}\ne\frac35.
\]
The contrast bias is $-21/80$. This is a fully stated population witness,
not a claim that the package necessarily clips or miscodes a user's model.
\citep{exec24_delicatessen_choice}

If $Y$ denotes a successful resolution, a stipulated 1,000-account
decision valued at 100 USD per additional success and costing 50,000 USD
has true net value 10,000 USD in that constructed population. Substituting
the clipped pseudo-target gives net value $-16,250$ USD. Both amounts are
conditional calculations, not bank returns. They show why the proposal
must preserve the estimand, nuisance specification and clipping sensitivity
through the financial decision. Threshold nonsmoothness also needs
derivative and uncertainty review; changing a threshold after inspecting
results is not a predeclared sensitivity design.
\citep{exec24_delicatessen_choice}

\subsection{Reference choice and the route to a qualified worker}
Direct adoption as a general causal service is deferred. A narrow isolated
integration remains a candidate after equation and sampling-law review,
environment locking and public known-law tests. Porting is also deferred:
it transfers responsibility for derivative, solver, shape and covariance
semantics without a measured maintenance benefit. The immediate decision
is to retain these selected sources and independent witnesses as a
reference. Unsupported certainty and automatic dependence correction are
rejected claims, not reasons to reject every scoped use of the package.
\citep{exec24_delicatessen_choice}

The scientific owner specifies identification, the target, independent
units and nuisance rates. The adapter owner records shapes, residuals,
rank, derivative settings, covariance normalization and refusal states.
The platform and rights owner resolves the complete environment and
redistribution obligations. Evaluation retains every request, including
nonconvergence, singularity, invalid domains, timeout and missing output.
For an authored 4,800 USD total over 240 requests, 200 finite results and
160 accepted results, costs are respectively 20, 24 and 30 USD per result
under those different denominators. No measured package economics is
inferred.~\citep{exec24_delicatessen_choice}

An agent may help translate a business question into a candidate equation
stack and propose unit, shape and residual checks. Human review must still
approve the causal assumptions, information clock and sampling law. CLI
and MCP responses should return the same typed target and uncertainty
status as a library call, including explicit refusal. Current inspected
sources do not establish arbitrary learner rates, unrestricted hidden-state
identification or universal dependent-data calibration. Those questions
need additional information or defensible restrictions. Execution of the package, independent validation and evaluation in a bank remain outstanding.
\citep{exec24_delicatessen_choice}


% SOURCE docs/release/2026-10-29/source/docs/writing-repair/2026-10-29/src/b3e5a34df36eb004.tex CHAR 0
\section{A method-specific reuse decision for Python DoubleML}
\label{exec23:doubleml-reuse-scope}
A uniform agent interface is most useful when its underlying scientific
operation has a precise target. Python DoubleML is a natural reference for
orthogonal-score estimation, but the package name does not identify one
estimand or one sampling law. The inspected source is pinned to commit
\texttt{6884b5f9816f180324d854a4d81b22e9d334f5e0}, dated 23 September 2026.
Five authenticated files cover the license, dependency declaration, data and
panel containers, and interactive-regression model. The checkout uses a dynamic
version scheme. It is not silently equated with the previously observed
installed version 0.11.4, whose selected local findings retain their original
scope. The R implementation, generic double-machine-learning methodology and
wrappers in other packages are separate comparison objects.
\citep{exec23_doubleml_reuse_source,sw_doubleml,runtime_doubleml_scope}

The inspected BSD three-clause notice requires retained notices and excludes
unauthorized endorsement; it is not a clearance of all transitive dependencies
or datasets. The source declares Python 3.10 or later and lower bounds on
scientific libraries, with a separate optional RD dependency. A named maintainer,
public issue tracker and a dated source revision provide concrete maintenance
identity, not a support-duration commitment. The proposed reuse route should
select an executable environment and a narrow method before comparing its
full cost with a new implementation.
\citep{exec23_doubleml_reuse_source}

\subsection{The score contract and its information clock}
\label{exec23:doubleml-reuse-score}
For a binary intervention $D$, observed outcome $Y$ and admissible pre-decision
covariates $X$, consider an average treatment effect under consistency,
conditional exchangeability and overlap. Let $g_d(x)$ denote the conditional
outcome mean in treatment group $d$ and $m(x)$ its treatment probability.
With cross-fitted or otherwise justified nuisance predictions, the unweighted
ATE pseudo-outcome is
\[
 Z_i=\widehat g_1(X_i)-\widehat g_0(X_i)
 +\frac{D_i\{Y_i-\widehat g_1(X_i)\}}{\widehat m(X_i)}
 -\frac{(1-D_i)\{Y_i-\widehat g_0(X_i)\}}{1-\widehat m(X_i)}.
\]
The inspected IRM implementation exposes separate ATE and ATTE weights,
normalization and propensity processing. Its ATE score has a negative
normalized-weight coefficient multiplying the target parameter and a residual
term with potentially distinct treatment-dependent weights. Therefore a weighted
or ATTE output must preserve its target and denominator, not merely inherit an
ATE label. The model rejects instrumental-variable data for this path and
requires exactly one zero-one treatment. Its external-prediction branch accepts
supplied nuisance arrays; their out-of-fold provenance remains a caller
obligation rather than an assumption established by their shape.
\citep{exec23_doubleml_reuse_source}

For an independent arithmetic reference, set $Y=(5,3,4,2)$,
$D=(1,0,1,0)$, $\widehat g_0=1$, $\widehat g_1=3$ and
$\widehat m=1/2$ for all four observations. The pseudo-outcomes are
$(6,-2,4,0)$, their mean is 2, and the stated empirical iid variance of the
mean is
\[
 \widehat V=\frac{1}{4^2}\sum_{i=1}^{4}(Z_i-2)^2=\frac52.
\]
This exact example tests score translation; it does not establish a
finite-sample confidence level or nuisance-learning performance. With an
entirely stipulated value of USD 10 per effect unit and cost USD 15, the
illustrative net is USD 5. In a separate one-residual example, replacing
propensity $0.001$ by $0.01$ changes a treated residual contribution of 1 from
1000 to 100. Clipping therefore matters to the numerical score and its bias
analysis; it cannot manufacture treatment support or identify an effect in a
population without overlap. No DoubleML code is run for these examples.
\citep{exec23_doubleml_reuse_reference}

The data container's default covariate selection includes otherwise unassigned
columns, and its multiple-treatment option can include other treatments as
covariates. A banking adapter should use an explicit covariate allowlist tied
to the decision's information clock. A post-servicing balance or a future
repayment outcome cannot become a valid adjustment variable merely because
it was present in the input table. The panel container has a separate static
mode; its default encodes groups by first treatment exposure, and it restricts
the input to one treatment column. In static mode it uses the unit identifier
as a cluster variable. Neither a first-adoption encoding nor the presence of
cluster labels makes the container an arbitrary longitudinal action history,
or proves an appropriate covariance law for shared calendar shocks.
\citep{exec23_doubleml_reuse_source}

\subsection{Reference now, a bounded adapter after qualification}
\label{exec23:doubleml-reuse-decision}
The proposal disposition is to retain Python DoubleML as a method-specific
reference and condition a future isolated IRM worker on explicit qualification.
Direct whole-package adoption is deferred because target, propensity processing,
score solving and uncertainty aggregation vary across families. A new language
port is not justified without a measured platform or cost advantage that
outweighs reproducing and maintaining those details. R DoubleML remains a
separate candidate for matched comparisons, and StatsPAI a separate integration
candidate. Neither receives automatic credit from the Python package's
source transparency or earlier installed-version results.
\citep{exec23_doubleml_reuse_source}

The conversion contract requires stable observation IDs and units, explicit
outcome and treatment definitions, nuisance-array ordering, saved folds and
training provenance, raw and processed propensities, normalization and weights.
Outputs should retain score components, the declared target, complete applicable
uncertainty and all failures. Before native qualification, inspect the selected
propensity processor, score solver, repeated-split aggregation and covariance
delegates at the same revision, freeze the environment and review its data and
dependency notices. Then compare a selected operation with an independent
reference on identical inputs, including refusals and uncertainty, under the
registered public or generated-data protocol. The scientific owner supplies
identification and sampling assumptions; the maintainer owns conversion,
versioning and regression checks; the operator owns custody and resource limits.
Changed source, target, sampling law or platform triggers re-review. This
source follow-up retains the earlier reference disposition. The recorded reuse
decision remains inconclusive because a complete release artifact, dependency
lock and original license inventory are still missing. No native integration
capability or general inference guarantee is awarded.


Across these comparisons, the most consequential differences are often changes
in the question: all customers versus treated customers, a marginal outcome
prediction versus treatment-specific regressions, or a row average versus an
average over customers. The retained numerical witnesses show how such changes
can reverse a financial recommendation even before estimation error is added.
The preferred reference routes therefore remain method-specific, with their
known restrictions and failed cases visible. If conditional exchangeability
cannot be defended, a more flexible learner does not complete the argument.
The next chapter examines what an instrument or eligibility threshold can
identify instead, and which population that information can support.

